\section{AIME I 2014}
        \begin{enumerate}
        	\item (\textit{AIME I 2014}) The 8 eyelets for the lace of a sneaker all lie on a rectangle, four equally spaced on each of the longer sides. The rectangle has a width of 50 mm and a length of 80 mm. There is one eyelet at each vertex of the rectangle. The lace itself must pass between the vertex eyelets along a width side of the rectangle and then crisscross between successive eyelets until it reaches the two eyelets at the other width side of the rectangle as shown. After passing through these final eyelets, each of the ends of the lace must extend at least 200 mm farther to allow a knot to be tied. Find the minimum length of the lace in millimeters. 


 
 \textbf{Answer: }790



\textbf{Solution}

The rectangle is divided into three smaller rectangles with a width of 50 mm and a length of $\dfrac{80}{3}$mm.  According to the Pythagorean Theorem (or by noticing the 8-15-17 Pythagorean triple), the diagonal of the rectangle is$\sqrt{50^2+\left(\frac{80}{3}\right)^2}=\frac{170}{3}$mm.  Since that on the lace, there are 6 of these diagonals, a width, and an extension of at least 200 mm on each side.  Therefore, the minimum of the lace in millimeters is
 \[6\times \dfrac{170}{3}+50+200\times 2=\boxed{790}.\]

 




	\item (\textit{AIME I 2014}) An urn contains $4$ green balls and $6$ blue balls. A second urn contains $16$ green balls and $N$ blue balls. A single ball is drawn at random from each urn. The probability that both balls are of the same color is $0.58$. Find $N$.


 
 \textbf{Answer: }144



\textbf{Solution}

First, we find the probability both are green, then the probability both are blue, and add the two probabilities. The sum should be equal to $0.58$. 

 The probability both are green is $\frac{4}{10}\cdot\frac{16}{16+N}$, and the probability both are blue is $\frac{6}{10}\cdot\frac{N}{16+N}$, so 
 \[\frac{4}{10}\cdot\frac{16}{16+N}+\frac{6}{10}\cdot\frac{N}{16+N}=\frac{29}{50}\]Solving this equation, 
 \[20\left(\frac{16}{16+N}\right)+30\left(\frac{N}{16+N}\right)=29\]Multiplying both sides by $16+N$, we get

 
 \begin{align*} 20\cdot16+30\cdot N&=29(16+N)\\ 320+30N&=464+29N\\ N&=\boxed{144} \end{align*}


 




	\item (\textit{AIME I 2014}) Find the number of rational numbers $r$, $0<r<1,$ such that when $r$ is written as a fraction in lowest terms, the numerator and the denominator have a sum of $1000$.


 
 \textbf{Answer: }200



\textbf{Solution 1 (Euclidean algorithm, fast)}

Let the numerator and denominator $x,y$ with $\gcd(x,y)=1$ and $x+y = 1000.$ Now if $\gcd(x,y) = 1$ then $\gcd(x,y) =\gcd(x,1000-x)= \gcd(x,1000-x-(-1)x)=\gcd(x,1000)=1.$ Therefore any pair that works satisfies $\gcd(x,1000)= 1.$ By Euler's totient theorem, there are $\phi(1000) = 400$ numbers relatively prime to 1000 from 1 to 1000. Recall that $r=\frac{x}{y}<1$ and note by Euclidean algorithm $\gcd(1000,1000-x)=1$, so we want $x<y=1000-x.$ Thus the $400$ relatively prime numbers can generate $\boxed{200}$ desired fractions.

 





\textbf{Solution 2}

If the initial manipulation is not obvious, consider the Euclidean Algorithm. Instead of using $\frac{n}{m}$ as the fraction to use the Euclidean Algorithm on, we can rewrite this as $\frac{500-x}{500+x}$ or 
 \[\gcd(500+x,500-x)=\gcd((500+x)+(500-x),500-x)=\gcd(1000,500-x).\] Thus, we want $\gcd(1000,500-x)=1$. You can either proceed as Solution $1$, or consider that no even numbers work, limiting us to $250$ choices of numbers and restricting $x$ to be odd. If $x$ is odd, $500-x$ is odd, so the only possible common factors $1000$ and $500-x$ can share are multiples of $5$. Thus, we want to avoid these. There are $50$ odd multiples of $5$ less than $500$, so the answer is $250-50=\boxed{200}$.

 





\textbf{Solution 3}

Say $r=\frac{d}{1000-d}$; then $1\leq d\leq499$. If this fraction is reducible, then the modulus of some number for $d$ is the same as the modulus for $1000-d$. Since $1000=2^3\cdot5^3$, that modulus can only be $2$ or $5$. This implies that if $d\mid2$ or $d\mid5$, the fraction is reducible. There are $249$ cases where $d\mid2$, $99$ where $d\mid5$, and $49$ where $d\mid(2\cdot5=10)$, so by PIE, the number of fails is $299$, so our answer is $\boxed{200}$.

 





\textbf{Solution 4}

We know that the numerator of the fraction cannot be even, because the denominator would also be even. We also know that the numerator cannot be a multiple of $5$ because the denominator would also be a multiple of $5$. Proceed by listing out all the other possible fractions and we realize that the numerator and denominator are always relatively prime. We have $499$ fractions to start with, and $250$ with odd numerators. Subtract $50$ to account for the multiples of $5$, and we get $\boxed{200}$.

 





\textbf{Solution 5}

We know that the set of these rational numbers is from $\dfrac{1}{999}$ to $\dfrac{499}{501}$ where each each element $\dfrac{n}{m}$ has $n+m =1000$ and $\dfrac{n}{m}$ is irreducible. 

 We note that $\dfrac{n}{m} =\dfrac{1000-m}{m}=\dfrac{1000}{m}-1$.Hence, $\dfrac{n}{m}$ is irreducible if $\dfrac{1000}{m}$ is irreducible, and $\dfrac{1000}{m}$ is irreducible if $m$ is not divisible by $2$ or $5$. Thus, the answer to the question is the number of integers between $501$ and $999$ inclusive that are not divisible by $2$ or $5$.

 We note there are $499$ numbers between $501$ and $999$, and

 
                \begin{itemize}
                    \item $249$ numbers are divisible by $2$\item $99$ numbers are divisible by $5$\item $49$ numbers are divisible by $10$
                \end{itemize}


                Using the Principle of Inclusion and Exclusion, we get that there are $499-249-99+49=200$ numbers between $501$ and $999$ are not divisible by either $2$ or $5$, so our answer is $\boxed{200}$.

 Euler's Totient Function can also be used to arrive at $400$ numbers relatively prime to $1000$, meaning $200$ possible fractions satisfying the necessary conditions. Solution 1 is a more detailed solution utilizing Euler's totient.

 





\textbf{Solution 6}

We notice that there are a total of $400$ fractions that are in simplest form where the numerator and denominator add up to $1000$. Because the numerator and denominator have to be relatively prime, there are $\varphi(1000)=400$ fractions. Half of these are greater than $1$, so the answer is $400\div2=\boxed{200}$

 - bedwarsnoob

 





\textbf{Solution 7}

Our fraction can be written in the form $\frac{1000 - a}{a} = \frac{1000}{a} - 1.$ Thus the fraction is reducible when $a$ divides $1000.$ We also want $500 < a < 1000.$ By \href{https://artofproblemsolving.com/wiki/index.php?title=PIE}{PIE}, the total values of $a$ that make the fraction reducible is, 
 \[249 + 99 - 49 = 299.\]By \href{https://artofproblemsolving.com/wiki/index.php?title=Complementary\_counting}{complementary counting}, the answer we want is $499 - 299 = \boxed{200}.$

 





\textbf{Solution 8 (Simplest)}

Suppose our fraction is $\frac{a}{b}$. The given condition means $a+b=1000$. Now, if $a$ and $b$ share a common factor greater than $1$, then the expression $a+b$ must also contain that common factor. This means our fraction cannot have a factor of $5$ or be even.

 There are $250$ fractions that aren’t even. From this, $50$ are divisible by $5$, which means the answer is $250-50=\boxed{200}$

 




	\item (\textit{AIME I 2014}) Jon and Steve ride their bicycles along a path that parallels two side-by-side train tracks running the east/west direction. Jon rides east at $20$ miles per hour, and Steve rides west at $20$ miles per hour. Two trains of equal length, traveling in opposite directions at constant but different speeds each pass the two riders. Each train takes exactly $1$ minute to go past Jon. The westbound train takes $10$ times as long as the eastbound train to go past Steve. The length of each train is $\frac{m}{n}$ miles, where $m$ and $n$ are relatively prime positive integers. Find $m+n$.


 
 \textbf{Answer: }049



\textbf{Solution 1}

For the purposes of this problem, we will use miles and minutes as our units; thus, the bikers travel at speeds of $\dfrac{1}{3}$ mi/min.

 Let $d$ be the length of the trains, $r_1$ be the speed of train 1 (the faster train), and $r_2$ be the speed of train 2.

 Consider the problem from the bikers' moving frame of reference. In order to pass Jon, the first train has to cover a distance equal to its own length, at a rate of $r_1 - \dfrac{1}{3}$. Similarly, the second train has to cover a distance equal to its own length, at a rate of $r_2 + \dfrac{1}{3}$. Since the times are equal and $d = rt$, we have that $\dfrac{d}{r_1 - \dfrac{1}{3}} = \dfrac{d}{r_2 + \dfrac{1}{3}}$. Solving for $r_1$ in terms of $r_2$, we get that $r_1 = r_2 + \dfrac{2}{3}$.

 Now, let's examine the times it takes the trains to pass Steve. This time, we augment train 1's speed by $\dfrac{1}{3}$

 




	\item (\textit{AIME I 2014}) Let the set $S = \{P_1, P_2, \dots, P_{12}\}$ consist of the twelve vertices of a regular $12$-gon. A subset $Q$ of $S$ is called communal if there is a circle such that all points of $Q$ are inside the circle, and all points of $S$ not in $Q$ are outside of the circle. How many communal subsets are there? (Note that the empty set is a communal subset.)


 
 \textbf{Answer: }134



\textbf{Solution}

By looking at the problem and drawing a few pictures, it quickly becomes obvious that one cannot draw a circle that covers $2$ disjoint areas of the $12$-gon without including all the vertices in between those areas. In other words, in order for a subset to be communal, all the vertices in the subset must be adjacent to one another. We now count the number of ways to select a row of adjacent vertices. We notice that for any subset size between $1$ and $11$, there are $12$ possible subsets like this (this is true because we can pick any of the $12$ vertices as a "starting" vertex, include some number of vertices counterclockwise from that vertex, and generate all possible configurations). However, we also have to include the set of all $12$ vertices, as well as the empty set. Thus, the total number is $12\cdot11 + 2 = \boxed{134}$.

 




	\item (\textit{AIME I 2014}) The graphs $y = 3(x-h)^2 + j$ and $y = 2(x-h)^2 + k$ have y-intercepts of $2013$ and $2014$, respectively, and each graph has two positive integer x-intercepts. Find $h$.


 
 \textbf{Answer: }036



\textbf{Solution 1}

Begin by setting $x$ to 0, then set both equations to $h^2=\frac{2013-j}{3}$ and $h^2=\frac{2014-k}{2}$, respectively. Notice that because the two parabolas have to have positive x-intercepts, $h\ge32$.

 We see that $h^2=\frac{2014-k}{2}$, so we now need to find a positive integer $h$ which has positive integer x-intercepts for both equations.

 Notice that if $k=2014-2h^2$ is -2 times a square number, then you have found a value of $h$ for which the second equation has positive x-intercepts. We guess and check $h=36$ to obtain $k=-578=-2(17^2)$.

 Following this, we check to make sure the first equation also has positive x-intercepts (which it does), so we can conclude the answer is $\boxed{036}$.

 
$\textbf{Can someone explain the logic of guessing 36?}$

 





\textbf{Solution 2}

Let $x=0$ and $y=2013$ for the first equation, resulting in $j=2013-3h^2$.  Substituting back in to the original equation, we get $y=3(x-h)^2+2013-3h^2$.

 Now we set $y$ equal to zero, since there are two distinct positive integer roots.  Rearranging, we get $2013=3h^2-3(x-h)^2$, which simplifies to $671=h^2-(x-h)^2$.  Applying difference of squares, we get $671=(2h-x)(x)$.

 Now, we know that $x$ and $h$ are both integers, so we can use the fact that $671=61\times11$, and set $2h-x=11$ and $x=61$ (note that letting $x=11$ gets the same result).  Therefore, $h=\boxed{036}$.

 Note that we did not use the second equation since we took advantage of the fact that AIME answers must be integers.  However, one can enter $h=36$ into the second equation to verify the validity of the answer. 

 Note on the previous note: we still must use the second equation since we could also use $671=671\times1$, yielding $h=336.$ This answer however does not check out with the second equation which is why it is invalid.

 





\textbf{Solution 3}

Similar to the first two solutions, we deduce that $\text{(-)}j$ and $\text{(-)}k$ are of the form $3a^2$ and $2b^2$, respectively, because the roots are integers and so is the $y$-intercept of both equations. So the $x$-intercepts should be integers also.

 The first parabola gives
 \[3h^2+j=3\left(h^2-a^2\right)=2013\]
 \[h^2-a^2=671\]And the second parabola gives
 \[2h^2+k=2\left(h^2-b^2\right)=2014\]
 \[h^2-b^2=1007\]

 We know that $671=11\cdot 61$ and that $1007=19\cdot 53$. It is just a fitting coincidence that the average of $11$ and $61$ is the same as the average of $19$ and $53$. That is $\boxed{036}$.

 To check, we have
 \[(h-a)(h+a)=671=11\cdot 61\]
 \[(h-b)(h+b)=1007=19\cdot 53\]Those are the only two prime factors of $671$ and $1007$, respectively. So we don't need any new factorizations for those numbers.

 $h+a=61,h-a=11\implies (h,a)=\{36,25\}$

 $h+b=53,h-b=19\implies (h,b)=\{36,17\}$

 Thus the common integer value for $h$ is $\boxed{036}$.

 
 \[y=3(x-h)^2+j\rightarrow y=3(x-11)(x-61)=3x^2-216x+2013\]
 \[y=2(x-h)^2+k\rightarrow y=2(x-19)(x-53)=2x^2-144x+2014\]

 





\textbf{Solution 4}

First, we expand both equations to get $y=3x^2-6hx+3h^2+j$ and $y=2x^2-4hx+2h^2+k$. The $y$-intercept for the first equation can be expressed as $3h^2+j$. From this, the x-intercepts for the first equation can be written as 

 
 \[x=h \pm \sqrt{(-6h)^2-4*3(3h^2+j)}=h \pm \sqrt{36h^2-12(2013)}=h \pm \sqrt{36h^2-24156}\]

 Since the $x$-intercepts must be integers, $\sqrt{36h^2-24156}$ must also be an integer. From solution 1, we know  $h$ must be greater than or equal to 32. We can substitute increasing integer values for $h$ starting from 32; we find that $h=36$. 

 We can test this result using the second equation, whose $x$-intercepts are 
 \[x=h \pm \sqrt{(-4h)^2-4*2(2h^2+k)}=h \pm \sqrt{16h^2-8(2014)}=h \pm \sqrt{16h^2-16112}\] Substituting 36 in for $h$, we get $h=36 \pm 68$, which satisfies the requirement that all x-intercepts must be (positive) integers.

 Thus, $h=\boxed{036}$.

 





\textbf{Solution 5}

We have the equation $y=3(x-h)^2 + j.$

 We know: $(x,y):(0,2013)$, so $h^2=2013/3 - j/3$ after plugging in the values and isolating $h^2$. Therefore, $h^2=671-j/3$.

 Lets call the x-intercepts $x_1$, $x_2$. Since both $x_1$ and $x_2$ are positive there is a relationship between $x_1$, $x_2$ and $h$. Namely, $x_1+x_2=2h$. The is because: $x_1-h=-(x_2-h)$, 

 Similarly, we know: $(x,y):(x_1,0)$, so $j=-3(x_1-h)^2$. Combining the two equations gives us 
 \[h^2=671+(x_1-h)^2\]
 \[h^2=671+x_1^2-2x_1h+h^2\]
 \[h=(671+x_1^2)/2x_1.\]

 Now since we have this relationship, $2h=x_1+x_2$, we can just multiply the last equation by 2(so that we get $2h$ on the left side) which gives us 
 \[2h=671/x_1+x_1^2/x^1\] 
 \[2h=671/x_1+x_1\] 
 \[x_1+x_2=671/x_1+x_1\] 
 \[x_2=671/x_1\] 
 \[x_1x_2=671.\] Prime factorization of 671 gives 11 and 61. So now we know $x_1=11$ and $x_2=61$. Lastly, we plug in the numbers,11 and 61, into $x_1+x_2=2h$, so $\boxed{h=36}$.

 





\textbf{Solution 6 (Vieta's solution)}

First, we start of exactly like solutions above and we find out that $j=2013-3h^2$ and $k=2014-2h^2$ We then plug j and k into $3(x-h)^2+j$ and $y=2(x-h)^2+k$ respectively. After that, we get two equations, $y=3x^2-6xh+2013$ and $y=2x^2-4xh+2014$. We can apply Vieta's. Let the roots of the first equation be $a, b$ and the roots of the second equation be $c, d$. Thus, we have that $a\cdot b=1007$, $a+b=2h$ and $c\cdot d=671$, $c+d=2h$. Simple evaluations finds that $\boxed{h=36}$

 





\textbf{Solution 7 (Copied from forums)}

Note that the resulting quadratics can be rewritten as $3x^2-6xh+2013$ and $2x^2-4xh+2014$ respectively. Therefore by Vieta's there exist roots $r_1$, $r_2$ of the former quadratic and $s_1$, $s_2$ of the latter such that
 \[r_1+r_2=s_1+s_2=2h,\,\,\, r_1r_2=671=61\times 11,\,\,\,\text{and}\,\,\, s_1s_2=1007=19\times 53.\]It so happens that $61+11=19+53=72$, so $2h=72\implies h=\boxed{036}$.

 





\textbf{Solution COPIED FROM ALCUMUS}

Not my solution, please take it down if it violates rules but this is copied from Alcumus:

 Setting $x=0$ in both equations, we get
 \[2013 = 3h^2 + j \quad \text{and} \quad 2014 = 2h^2 + k.\]Solving for $j$ and $k,$ we can rewrite the given equations as
 \[y = 3(x-h)^2 + (2013-3h^2) \quad \text{and} \quad y = 2(x-h)^2 + (2014-2h^2),\]or
 \[y = 3x^2 - 6xh + 2013 = 3(x^2-2hx+671) \quad \text{ and } \quad y = 2x^2 - 4hx + 2014 = 2(x^2 - 2hx + 1007).\]The left equation has positive integer roots, which must multiply to $671$ and sum to $2h.$ Similarly, the right equation has positive integer roots, which must multiply to $1007$ and sum to $2h.$ Since $671 = 61 \cdot 11$ and $1007 = 19 \cdot 53,$ we see that
 \[2h = 61 + 11 = 19 + 53 = 72,\]so $h = \boxed{36}.$

 




	\item (\textit{AIME I 2014}) Let $w$ and $z$ be complex numbers such that $|w| = 1$ and $|z| = 10$. Let $\theta = \arg \left(\frac{w-z}{z}\right)$. The maximum possible value of $\tan^2 \theta$ can be written as $\frac{p}{q}$, where $p$ and $q$ are relatively prime positive integers. Find $p+q$. (Note that $\arg(w)$, for $w \neq 0$, denotes the measure of the angle that the ray from $0$ to $w$ makes with the positive real axis in the complex plane.)


 
 \textbf{Answer: }100



\textbf{Solution 1}

Let $w = \operatorname{cis}{(\alpha)}$ and $z = 10\operatorname{cis}{(\beta)}$. Then, $\dfrac{w - z}{z} = \dfrac{\operatorname{cis}{(\alpha)} - 10\operatorname{cis}{(\beta)}}{10\operatorname{cis}{\beta}}$.

 Multiplying both the numerator and denominator of this fraction by $\operatorname{cis}{(-\beta)}$ gives us:

 $\dfrac{w - z}{z} = \dfrac{1}{10}\operatorname{cis}{(\alpha - \beta)} - 1 = \dfrac{1}{10}\cos{(\alpha - \beta)} + \dfrac{1}{10}i\sin{(\alpha - \beta)} - 1$.

 We know that $\tan{\theta}$ is equal to the imaginary part of the above expression divided by the real part. Let $x = \alpha - \beta$. Then, we have that:

 $\tan{\theta} = \dfrac{\sin{x}}{\cos{x} - 10}.$

 We need to find a maximum of this expression, so we take the derivative:

 Note (not from author): To take the derivative, we need to use the \href{https://en.wikipedia.org/wiki/Quotient\_rule}{Quotient Rule}. In this case, 
 \[\frac{d}{dx}\left(\frac{\sin x}{\cos x-10}\right)=\frac{\cos x(\cos x-10)-(-\sin x)\sin x}{(\cos x-10)^2}=\dfrac{1 - 10\cos{x}}{(\cos{x} - 10)^2}\]

 Thus, we see that the maximum occurs when $\cos{x} = \dfrac{1}{10}$. Therefore, $\sin{x} = \pm\dfrac{\sqrt{99}}{10}$, and $\tan{\theta} = \pm\dfrac{\sqrt{99}}{99}$. Thus, the maximum value of $\tan^2{\theta}$ is $\dfrac{99}{99^2}$, or $\dfrac{1}{99}$, and our answer is $1 + 99 = \boxed{100}$.

 





\textbf{Solution 2}

Without the loss of generality, one can let $z$ lie on the positive x axis and since $arg(\theta)$ is a measure of the angle if $z=10$ then $arg(\dfrac{w-z}{z})=arg(w-z)$ and we can see that the question is equivalent to having a triangle $OAB$ with sides $OA =10$ $AB=1$ and $OB=t$ and trying to maximize the angle $BOA$
                    \begin{figure}[H]
                        \centering
                        \includegraphics[width=0.5\textwidth]{Solution/AIME_I_2014_7_Solution_Figure_1.png}
                    \end{figure}
                    

 using the Law of Cosines we get:$1^2=10^2+t^2-t*10*2\cos\theta$rearranging:
 \[20t\cos\theta=t^2+99\]solving for $\cos\theta$ we get:

 
 \[\frac{99}{20t}+\frac{t}{20}=\cos\theta\]if we want to maximize $\theta$ we need to minimize $\cos\theta$, using AM-GM inequality we get that the minimum value for $\cos\theta= 2\left(\sqrt{\dfrac{99}{20t}\cdot\dfrac{t}{20}}\right)=2\sqrt{\dfrac{99}{400}}=\dfrac{\sqrt{99}}{10}$hence using the identity $\tan^2\theta=\sec^2\theta-1$we get $\tan^2\theta=\frac{1}{99}$and our answer is $1 + 99 = \boxed{100}$.

 Note: You can also realize that the max $\theta$ is when the line from $0$ is tangent to the circle of radius $1$ centered at $10.$

 





\textbf{Solution 3}

Note that $\frac{w-z}{z}=\frac{w}{z}-1$, and that $\left|\frac{w}{z}\right|=\frac{1}{10}$. Thus $\frac{w}{z}-1$ is a complex number on the circle with radius $\frac{1}{10}$ and centered at $-1$ on the complex plane. Let $\omega$ denote this circle. 

 Let $A$ and $C$ be the points that represent $\frac{w}{z}-1$ and $-1$ respectively on the complex plane. Let $O$ be the origin. In order to maximize $\tan^2(\theta)$, we need to maximize $\angle{AOC}$. This angle is maximized when $AO$ is tangent to $\omega$. Using the Pythagorean Theorem, we get 

 
 \[AO^2=1^2-\left(\frac{1}{10}\right)^2=\frac{99}{100}\]

 Thus

 
 \[\tan^2(\theta)=\frac{AC^2}{AO^2}=\frac{1/100}{99/100}=\frac{1}{99}\]

 And the answer is $1+99=\boxed{100}$.

 





\textbf{Solution 4}

Let $\omega = a+bi$ and $z = c+di$. Note 
 \[(a^2+b^2)(c^2+d^2) =(ac+bd)^2 + (ad-bc)^2 =  100\] Further, 
 \[\theta = \dfrac{ac+bd - (c^2+d^2) + (-ad + bc)i}{100}\] Let $ac+bd = x$. Note $|x| \leq 10$. Then 
 \[\tan^2{\theta}= \frac{(-ad+bc)^2}{(ac+bd-100)^2} = \frac{100-x^2}{(100-x)^2}\] Taking the derivative and setting equal to zero (or via substitution and quadratic simplification), we obtain max occurs at $x = 1$. Substituting back in gives $\tan^2\theta = \frac{1}{99}$, and an answer of $1+99=\boxed{100}$.

 





\textbf{Solution 5}

Let $\omega=e^{i\alpha}, z=10e^{i\beta}$. We have\begin{align*}\frac{\omega}{z}-1&=\frac{e^{i\alpha}}{10e^{i\beta}}-1\\&=\frac{1}{10}e^{i(\alpha-\beta)}-1\\&=\left(\frac{1}{10}\cos(\alpha-\beta)-1\right)+\frac{1}{10}i\sin{(\alpha-\beta)}\\\end{align*}
Let $x=\alpha-\beta$, then $\tan\theta$ evaluates to
 \[\tan\theta=\frac{\frac{1}{10}\sin{x}}{\frac{1}{10}\cos{x}-1}=\frac{\sin{x}}{\cos{x}-10}\]We introduce an elementary method to solve the bound of this expression. We set it equal a new variable:
 \[\frac{\sin{x}}{\cos{x}-10}=k\]We are interested in the maximal of $k$, continuing manipulating:\begin{align*}\frac{\sin{x}}{\cos{x}-10}&=k\\\sin{x}&=k\cdot\cos{x}-10k\\k\cdot\cos{x}-\sin{x}&=10k\end{align*}
By the angle-addition formula, let $\tan\phi=\frac{1}{k}$, and we note that\begin{align*}k\cdot\cos{x}-\sin{x}&=\sqrt{1+k^2}\left(\frac{k}{\sqrt{1+k^2}}\cos{x}-\frac{1}{k}\sin{x}\right)\\&=\sqrt{1+k^2}\Bigl(\cos{x}\cdot\cos{\phi}-\sin{x}\cdot\sin{\phi}\Bigr)\\&=\sqrt{1+k^2}\cos{(x+\phi)}\end{align*}
Thus, $10k=\sqrt{1+k^2}\cos{(x+\phi)}$. It follows that
 \[-1\le\frac{10k}{\sqrt{1+k^2}}\le 1\]Solving yields $-\frac{1}{\sqrt{99}}\le k \le\frac{1}{\sqrt{99}}$. Therefore, the maximal of $\tan\theta$ evaluates to
 \[\tan^{2}\theta=\max{(\lvert k\rvert)}^2=\frac{1}{99}\]The requested sum is $1+99=\boxed{100}$.

 Note: If you are unfamiliar with the method above, note that we can also obtain the bound of $k\cdot \cos x-\sin x$ by the Cauchy-Schwarz inequality:
 \[(k\cdot \cos x-\sin x)^2\le (k^2+1)(\cos^2 x+\sin^2 x)\iff  -\sqrt{k^2+1}\le k\cdot \cos x-\sin x\le\sqrt{k^2+1}.\]

 


 




	\item (\textit{AIME I 2014}) The positive integers $N$ and $N^2$ both end in the same sequence of four digits $abcd$ when written in base 10, where digit $a$ is not zero. Find the three-digit number $abc$.


 
 \textbf{Answer: }937



\textbf{Solution 1 (similar to Solution 3)}

We have that $N^2 - N = N(N - 1)\equiv 0\mod{10000}$

 Thus, $N(N-1)$ must be divisible by both $5^4$ and $2^4$. Note, however, that if either $N$ or $N-1$ has both a $5$ and a $2$ in its factorization, the other must end in either $1$ or $9$, which is impossible for a number that is divisible by either $2$ or $5$. Thus, one of them is divisible by $2^4 = 16$, and the other is divisible by $5^4 = 625$. Noting that $625 \equiv 1\mod{16}$, we see that $625$ would work for $N$, except the thousands digit is $0$. The other possibility is that $N$ is a multiple of $16$ and $N-1$ is a multiple of $625$. In order for this to happen, 
 \[N-1 \equiv -1 \pmod {16}.\] Since $625 \equiv 1 \pmod{16}$, we know that $15 \cdot 625 = 9375 \equiv 15 \equiv -1 \mod{16}$. Thus, $N-1 = 9375$, so $N = 9376$, and our answer is $\boxed{937}$.

 





\textbf{Solution 2 (bashing)}

let $N= 10000t+1000a+100b+10c+d$ for positive integer values $t,a,b,c,d$.When we square $N$ we get that 
 \begin{align*} N^2 &=(10000t+1000a+100b+10c+d)^2\\ &=10^8t^2+10^6a^2+10^4b^2+10^2c^2+d^2+2(10^7ta+10^6tb+10^5tc+10^4td+10^5ab+10^4ac+10^3bc+10^ad+10^2bd+10cd) \end{align*}


 However, we don't have to deal with this whole expression but only with its last 4 digits so it is suffices to consider only:
 \[2000ad+2000bc+100c^2+200bd+20cd+d^2.\]Now we need to compare each decimal digit with $1000a+100b+10c+d$ and see whether the digits are congruent in base 10.we first consider the ones digits:

 $d^2\equiv d \pmod{10}.$

 This can happen for only 3 values : 1, 5 and 6.

 We can try to solve each case:

 
                \begin{itemize}
                    \item Case 1 $(d=1)$
                \end{itemize}


                Considering the tenths place, we have that:

 $20cd=20c\equiv 10c \pmod {100}$so $c= 0$.

 Considering the hundreds place we have that 

 $200bd+100c^2= 200b \equiv 100b \pmod{1000}$so again $b=0$

 now considering the thousands place we have that 

 $2000ad+2000bc = 2000a \equiv 1000a \pmod {10000}$so we get $a=0$ but $a$ cannot be equal to $0$ so we consider $d=5.$

 
                \begin{itemize}
                    \item Case 2 $(d=5)$
                \end{itemize}


                considering the tenths placewe have that:

 $20cd+20=100c+20\equiv 20 \equiv 10c \mod  {100}$( the extra $20$ is carried from $d^2$ which is equal to $25$)so $c=2$

 considering the hundreds place we have that 

 $200bd+100c^2+100c= 1000b+600 \equiv600\equiv 100b \pmod{1000}$( the extra $100c$ is carried from the tenths place)so $b=6$

 now considering the thousands place we have that 

 $2000ad+2000bc +1000b= 10000a+24000+ 6000\equiv0\equiv 1000a \pmod {10000}$( the extra $1000b$ is carried from the hundreds place)so a is equal 0 again

 
                \begin{itemize}
                    \item Case 3$(d=6)$
                \end{itemize}


                considering the tenths placewe have that:

 $20cd+30=120c+30\equiv 30+20c \equiv 10c \pmod {100}$( the extra $20$ is carried from $d^2$ which is equal to $25$)if $c=7$ then we have 

 $30+20 \cdot 7 \equiv 70\equiv7 \cdot 10 \pmod{100}$

 so $c=7$

 considering the hundreds place we have that 

 $200bd+100c^2+100c+100= 1200b+4900+800 \equiv200b+700\equiv 100b \pmod{1000}$( the extra $100c+100$ is carried from the tenths place)

 if $b=3$ then we have 

 $700+200 \cdot 3 \equiv 300\equiv3 \cdot 100 \pmod {1000}$

 so $b=3$

 now considering the thousands place we have that 

 $2000ad+2000bc +1000b+5000+1000= 12000a+42000+ 3000+6000\equiv0\equiv 2000a+1000\equiv 1000a \pmod {10000}$( the extra $1000b+6000$ is carried from the hundreds place)

 if $a=9$ then we have 

 $2000 \cdot 9+1000 \equiv 9000\equiv9 \cdot 1000 \pmod {1000}$

 so $a=9$

 so we have that the last 4 digits of $N$ are $9376$and  $abc$ is equal to $\boxed{937}$

 





\textbf{Solution 3 (general)}

By the Chinese Remainder Theorem, the equation $N(N-1)\equiv 0\pmod{10000}$ is equivalent to the two equations:
 \begin{align*} N(N-1)&\equiv 0\pmod{16},\\ N(N-1)&\equiv 0\pmod{625}. \end{align*}
Since $N$ and $N-1$ are coprime, the only solutions are when $(N\mod{16},N\mod{625})\in\{(0,0),(0,1),(1,0),(1,1)\}$.

 Let 
 \[\varphi:\mathbb Z/10000\mathbb Z\to\mathbb Z/16\mathbb Z\times\mathbb Z/625\mathbb Z,\] 
 \[x\mapsto (x\mod{16},x\mod{625}).\] The statement of the Chinese Remainder theorem is that $\varphi$ is an isomorphism between the two rings. In this language, the solutions are $\varphi^{-1}(0,0)$, $\varphi^{-1}(0,1)$, $\varphi^{-1}(1,0)$, and $\varphi^{-1}(1,1)$. Now we easily see that 
 \[\varphi^{-1}(0,0)=0\] and 
 \[\varphi^{-1}(1,1)=1.\] Noting that $625\equiv 1\pmod{16}$, it follows that 
 \[\varphi^{-1}(1,0)=625.\] To compute $\varphi^{-1}(0,1)$, note that 
 \[(0,1)=15(1,0)+(1,1)\] in 
 \[\mathbb Z/16\mathbb Z\times\mathbb Z/625\mathbb Z,\] so since $\varphi^{-1}$ is linear in its arguments (by virtue of being an isomorphism), 
 \[\varphi^{-1}(0,1)=15\varphi^{-1}(1,0)+\varphi^{-1}(1,1)=15\times 625+1=9376.\]

 The four candidate digit strings $abcd$ are then $0000,0001,0625,9376$. Of those, only $9376$ has nonzero first digit, and therefore the answer is $\boxed{937}$.

 





\textbf{Solution 4 (semi-bashing)}


                \begin{itemize}
                    \item Note - $\overline{abcd}$ means the number formed when the digits represented by $a$, $b$, $c$, and $d$ are substituted in. $\overline{abcd}\ne a\times b\times c\times d$.
                \end{itemize}


                WLOG, we can assume that $N$ is a 4-digit integer $\overline{abcd}$. Note that the only $d$ that will satisfy $N$ will also satisfy $d^2\equiv d\pmod{10}$, as the units digit of $\overline{abcd}^2$ is affected only by $d$, regardless of $a$, $b$, or $c$.

 By checking the numbers 0-9, we see that the only possible values of $d$ are $d=0, 1, 5, 6$.

 Now, we seek to find $c$. Note that the only $\overline{cd}$ that will satisfy $N$ will also satisfy $\overline{cd}^2 \equiv \overline{cd}\pmod{100}$, by the same reasoning as above - the last two digits of $\overline{abcd}^2$ are only affected by $c$ and $d$. As we already have narrowed choices for $d$, we start reasoning out.

 First, we note that if $d=0$, then $c=0$, as a number ending in 0, and therefore divisible by 10, is squared, the result is divisible by 100, meaning it ends in two 0's. However, if $N$ ends in $00$, then recursively, $a$ and $b$ must be $0$, as a number divisible by 100 squared ends in four zeros. As $a$ cannot be 0, we throw out this possibility, as the only solution in this case is $0$.

 Now, let's assume that $d=1$. $\overline{cd}$ is equal to $10c + d = 10c + 1$. Squaring this gives $100c^2 + 20c + 1$, and when modulo 100 is taken, it must equal $10c + 1$. As $c$ is an integer, $100c^2$ must be divisible by 100, so $100c^2+20c+1 \equiv 20c + 1\pmod{100}$, which must be equivalent to $10c + 1$. Note that this is really $\overline{(2c)1}$ and $\overline{c1}$, and comparing the 10's digits. So really, we're just looking for when the units digit of $2c$ and $c$ are equal, and a quick check reveals that this is only true when $c=0$.However, if we extend this process to find $b$ and $a$, we'd find that they are also 0. The only solution in this case is $1$, and since $a=0$ here, this is not our solution. Therefore, there are no valid solutions in this case.

 Let's assume that $d=5$. Note that $(10c + 5)^2 = 100c^2 + 100c + 25$, and when modulo $100$ is taken, $25$ is the remainder. So all cases here have squares that end in 25, so $\overline{cd}=25$ is our only case here. A quick check reveals that $25^2=625$, which works for now.

 Now, let $d=6$. Note that $(10c + 6)^2 = 100c^2 + 120c + 36$. Taking modulo 100, this reduces to $20c+36$, which must be equivalent to $10c+6$. Again, this is similar to $\overline{(2c+3)6}$ and $\overline{c6}$, so we see when the units digits of $2c+3$ and $c$ are equal. To make checking faster, note that $2c$ is necessarily even, so $2c+3$ is necessarily odd, so $c$ must be odd. Checking all the odds reveals that only $c=3$ works, so this case gives $76$. Checking quickly $76^2 = 5776$, which works for now.

 Now, we find $b$, given two possibilities for $\overline{cd}$. 

 Start with $\overline{cd} = 25$. $\overline{bcd} = 100b + \overline{cd} = 100b + 25.$ Note that if we square this, we get $10000b^2 + 5000b + 625$, which should be equivalent to $100b + 25$ modulo 1000. Note that, since $b$ is an integer, $10000b^2 + 5000 + 625$ simplifies modulo 1000 to $625$. Therefore, the only $\overline{bcd}$ that works here is $625$. $625^2 = 390625$.

 Now, assume that $\overline{cd}=76$. We have $100b + 76$, and when squared, becomes $10000b^2 + 15200b + 5776$, which, modulo 1000, should be equivalent to $100b+76$. Reducing $10000b^2 + 15200b + 5776$ modulo $1000$ gives $200b + 776$. Using the same technique as before, we must equate the hundreds digit of $\overline{(2b+7)76}$ to $\overline{b76}$, or equate the units digit of $2b+7$ and $b$. Since $2b+7$ is necessarily odd, any possible $b$'s must be odd. A quick check reveals that $b=3$ is the only solution, so we get a solution of $376$. $376^2 = 141376$.

 Finally, we solve for $a$. Start with $\overline{bcd}=625$. We have $1000a + 625$, which, squared, gives 
 \[1000000a^2 + 1250000a + 390625,\] and reducing modulo 10000 gives simply 625. So $\overline{abcd}=625$. However, that makes $a=0$. Therefore, no solutions exist in this case.

 We turn to our last case, $\overline{bcd}=376$. We have 
 \[1000a + 376^2 = 1000000a^2 + 752000a + 141376,\] and reducing modulo $10000$ gives $2000a + 1376$, which must be equivalent to $1000a + 376$. So we must have $\overline{(2a+1)376}$ being equivalent to $\overline{a376}$ modulo 1000. So, the units digit of $2a+1$ must be equal to $a$. Since $2a+1$ is odd, $a$ must be odd. Lo and behold, the only possibility for $a$ is $a=3$. Therefore, $\overline{abcd}=9376$, so our answer is $\boxed{937}$.

 





\textbf{Solution 5 (easy observations)}

We are given that $abcd^2$ ends in $abcd$, which means that $d^2$ ends in $d$. The only possible values for which this works is $d=0,1,5,6$. Notice that if we set $d=0$, we will have $c=0, b=0, a=0$, which is not valid since $a\neq0$. Furthermore, we can see that $d\neq1$ by quickly going over $11^2, 12^2, ..., 19^2$ and seeing that none of those numbers work. From number sense, we know that $B25^2$ ends in $...625$ regardless of $B$. So, in this case, $b=6$. However, notice that $76^2$ ends in $...76$, and we can quickly find that $376^2=141376$. Going back to $A625^2$, we notice that no matter what $A$ is, $A625^2=...0625$. Since $A\neq0$, this case is not valid (write out $A625^2$ and multiply to see that the thousands digit is $5a+3+2+5+5a=10a+10$ which always ends in $0$). We turn our focus to $A376^2$. Write out $376^2$ and multiply it out to see that the thousands digit is $6a+2+6+2+6a+1=12a+1$. There is an extra $1$ due to carry over. Quickly testing numbers, the only possible value for $A$ is $9$. So, $abcd=9376$ and our answer is $abc=\boxed{937}$.

 





\textbf{Solution 6 (Exact same as 3 but with less fancy words)}

The problem statement is equivalent to $N^2 \equiv N\pmod{10000}$, and subtracting $N$ from both sides and factoring yields $N(N-1) \equiv 0\pmod{10000}$. By the Chinese Remainder Theorem, the equation $N(N-1)\equiv 0\pmod{10000}$ is equivalent to the two equations:
 \begin{align*} N(N-1)&\equiv 0\pmod{16},\\ N(N-1)&\equiv 0\pmod{625}. \end{align*}


 So, 
 \begin{align*} N&\equiv 0, 1\pmod{16},\\ N&\equiv 0, 1\pmod{625}. \end{align*}


 If $N\equiv 0\pmod{16}$ and $N\equiv 0\pmod{625}$, then $N\equiv 0\pmod{10000}$, but this contradicts $a \neq 0$.

 If $N\equiv 1\pmod{16}$ and $N\equiv 1\pmod{625}$, then $N\equiv 1\pmod{10000}$, but this contradicts $a \neq 0$.

 
If $N\equiv 1\pmod{16}$ and $N\equiv 0\pmod{625}$, then $N\equiv 625\pmod{10000}$, but this contradicts $a \neq 0$.

 If $N\equiv 0\pmod{16}$ and $N\equiv 1\pmod{625}$, then $N\equiv 9376\pmod{10000}$. We can square this to check: $9376^2=87909376$. Our answer is $\boxed{937}$.

 




	\item (\textit{AIME I 2014}) Let $x_1< x_2 < x_3$ be the three real roots of the equation $\sqrt{2014} x^3 - 4029x^2 + 2 = 0$. Find $x_2(x_1+x_3)$.


 
 \textbf{Answer: }002



\textbf{Solution 1}

Substituting $n$ for $2014$, we get 
 \[\sqrt{n}x^3 - (1+2n)x^2 + 2 = \sqrt{n}x^3 - x^2 - 2nx^2 + 2\] 
 \[= x^2(\sqrt{n}x - 1) - 2(nx^2 - 1) = 0\] Noting that $nx^2 - 1$ factors as a difference of squares to 
 \[(\sqrt{n}x - 1)(\sqrt{n}x+1)\] we can factor the left side as 
 \[(\sqrt{n}x - 1)(x^2 - 2(\sqrt{n}x+1))\] This means that $\frac{1}{\sqrt{n}}$ is a root, and the other two roots are the roots of $x^2 - 2\sqrt{n}x - 2$. Note that the constant term of the quadratic is negative, so one of the two roots is positive and the other is negative. In addition, by Vieta's Formulas, the roots sum to $2\sqrt{n}$, so the positive root must be greater than $2\sqrt{n}$ in order to produce this sum when added to a negative value. Since $0 < \frac{1}{\sqrt{2014}} < 2\sqrt{2014}$ is clearly true, $x_2 = \frac{1}{\sqrt{2014}}$ and $x_1 + x_3 = 2\sqrt{2014}$. Multiplying these values together, we find that $x_2(x_1+x_3) = \boxed{002}$.

 





\textbf{Solution 2}

From Vieta's formula, we know that 
 \[x_1x_2x_3 = \dfrac{-2}{\sqrt{2014}}\] 
 \[x_1 + x_2 + x_3 = \dfrac{4029}{\sqrt{2014}}\] and 
 \[x_1x_2 + x_2x_3 + x_1x_3 = 0\] Thus, we know that 
 \[x_2(x_1 + x_3) = -x_1x_3\]

 Now consider the polynomial with roots $x_1x_2, x_2x_3,$ and $x_1x_3$. Expanding the polynomial 
 \[(x - x_1x_2)(x - x_2x_3)(x - x_1x_3)\]we get the polynomial 
 \[x^3 - (x_1x_2 + x_2x_3 + x_1x_3)x^2 + (x_1x_2x_3)(x_1 + x_2 + x_3)x - (x_1x_2x_3)^2\] Substituting the values obtained from Vieta's formulae, we find that this polynomial is 
 \[x^3 - \dfrac{8058}{2014}x - \dfrac{4}{2014}\] We know $x_1x_3$ is a root of this polynomial, so we set it equal to 0 and simplify the resulting expression to 
 \[1007x^3 - 4029x - 2 = 0\]

 Given the problem conditions, we know there must be at least 1 integer solution, and that it can't be very large (because the $x^3$ term quickly gets much larger/smaller than the other 2). Trying out some numbers, we quickly find that $x = -2$ is a solution. Factoring it out, we get that 
 \[1007x^3 + 4029x - 2 = (x+2)(1007x^2 - 2014x - 1)\] Since the other quadratic factor clearly does not have any integer solutions and since the AIME has only positive integer answers, we know that this must be the answer they are looking for. Thus, 
 \[x_1x_3 = -2\] so 
 \[-x_1x_3 = \boxed{002}\]and we're done.

 





\textbf{Solution 3}

Observing the equation, we notice that the coefficient for the middle term $-4029$ is equal to 
 \[-2{\sqrt{2014}}^2-1\]. 

 Also notice that the coefficient for the ${x^3}$ term is $\sqrt{2014}$. Therefore, if the original expression was to be factored into a linear binomial and a quadratic trinomial, the $x$ term of the binomial would have a coefficient of $\sqrt{2014}$. Similarly, the $x$ term of the trinomial would also have a coefficient of $\sqrt{2014}$. The factored form of the expression would look something like the following:
 \[({\sqrt{2014}}x-a)(x^2-n{\sqrt{2014}}x-b)\] where ${a, b,c}$ are all positive integers (because the ${x^2}$ term of the original expression is negative, and the constant term is positive), and 
 \[{ab=2}\]

 Multiplying this expression out gives 
 \[{{\sqrt{2014}x^3-(2014n+a)x^2+(an{\sqrt{2014}}-b{\sqrt{2014}})x+ab}}\] Equating this with the original expression gives 
 \[{2014n+a}=-4029\] The only positive integer solutions of this expression is $(n, a)=(1, 2015)$ or $(2, 1)$. If $(n, a)=(1, 2015)$ then setting ${an{\sqrt{2014}}-b{\sqrt{2014}}}=0$ yields ${b=2015}$ and therefore ${ab=2015^2}$ which clearly isn't equal to $2$ as the constant term. Therefore, $(n, a)=(2, 1)$ and the factored form of the expression is:
 \[({\sqrt{2014}}x-1)(x^2-2{\sqrt{2014}}x-2)\] Therefore, one of the three roots of the original expression is 
 \[{x=\dfrac{1}{\sqrt{2014}}}\]Using the quadratic formula yields the other two roots as 
 \[{x={\sqrt{2014}}+{\sqrt{2016}}}\] and 
 \[{x={\sqrt{2014}}-{\sqrt{2016}}}\] Arranging the roots in ascending order (in the order $x_1<x_2<x_3$), 
 \[{\sqrt{2014}}-{\sqrt{2016}}<\dfrac{1}{\sqrt{2014}}<{\sqrt{2014}}+{\sqrt{2016}}\]Therefore, 
 \[x_2(x_1+x_3)=\dfrac{1}{\sqrt{2014}}{2\sqrt{2014}}=\boxed{002}\]

 





\textbf{Solution 4}

By Vieta's, we are seeking to find $x_2(x_1+x_3)=x_1x_2+x_2x_3=-x_1x_3=\frac{2}{\sqrt{2014}x_2}$.

 Substitute $n=-x_1x_3$ and $x_2=\frac{2}{\sqrt{2014}n}$. 

 Substituting this back into the original equation, we have $\frac{4}{1007n^3}-\frac{8058}{1007n^2}+2=0$, so $2n^3-\frac{8058}{1007}n+\frac{4}{1007}=2n^3-\frac{8058n-4}{1007}=0$. 

 Hence, $8058n-4\equiv 2n-4 \equiv 0 \pmod{1007}$, and so $n\equiv 2\pmod{1007}$. 

 But since $n\le 999$ because it is our desired answer, the only possible value for $n$ is $\boxed{002}$

 Stormersyle & mathleticguyyy

 





\textbf{Solution 5}

Let $x =\frac{y}{\sqrt{2014}}.$

 The original equation simplifies to $\frac{y^3}{2014} -\frac{4029y^2}{2014}+2 = 0 \implies y^3 - 4029y^2 + 4028=0.$ 

 Here we clearly see that $y=1$ is a root. 

 Dividing $y-1$ from the sum we find that $(y-1)(y^2-4028y-4028)=0.$ From simple bounding we see that $y=1$ is the middle root. 

 Therefore $x_{2}(x_{1}+x_{3}) =\frac{1}{\sqrt{2014}} \cdot\frac{4028}{\sqrt{2014}} = \boxed{002}.$

 





\textbf{Solution 6}

$a=\sqrt{2014}$ looks arbitrary and related to $b=-4019$ ($b=-(2a^2+1)$), so let k = $\sqrt{2014}$.

 The equation becomes $0 = kx^3 - (2k^2 + 1)x^2 + 2$. Since we want to relate k and x, we should solve for one of them. We can't solve for x, since that would require the cubic formula, so we solve for k, and express it in terms of a quadratic, and apply the quadratic formula.

 We get the roots are:

 $k = \frac{1}{x}$, and $k = \frac{x}{2} - \frac{1}{x}$.

 In the first case, $x = \frac{1}{k} = \frac{1}{\sqrt{2014}}$.

 In the second case, $x^2 - 2\sqrt{2014} - 2 = 0$. The solutions are $\sqrt{2014} \pm \sqrt{2016}$. The sum of these 2 solutions is $2 \sqrt{2014}$, and $\frac{1}{\sqrt{2014}}$ is the middle solution, and thus, $(x_1 + x_3) \cdot x_2 = 2$

 





\textbf{Solution 7 (Estimation and Intuitive Function Analysis)}

We will estimate the roots of the polynomial. (This strategy normally doesn't work on AIME \#9, but playing around with a function is often good strategy for getting an intuition for the problem. In this problem, estimation happens to be a valid solution path. It isn't a proof, but given the constraint that the answer is an integer, we can be certain that our answer is correct.)

 
Let $p(x) = \sqrt{2014}x^3-4029x^2+2=0$. We start by estimating $p(-1)$, $p(0)$, and $p(1)$ (A natural first step for function analysis.):

 
$p(-1)\approx -45-4029+2 \approx -4000$

 $p(0) = 2$

 $p(1) \approx 45-4029+2 \approx -4000$

 
We conclude by Intermediate Value Theorem (or just common sense), that there is a root on $(-1, 0)$ and another root on $(0, 1)$.

 
We know that $p(1) < 0$ and that $\lim_{x\to\infty} p(x) = \infty$. We conclude that the third root is on $(1, \infty)$. Therefore, $x_1 \in (-1, 0)$, $x_2 \in (0, 1)$, and $x_3 \in (1, \infty)$.

 
We will estimate $x_3$. For $x > 1$, the constant term of $p(x)$ is negligible. We simplify and get $p(x_3) \approx \sqrt{2014}x_{3}^3 - 4029x_{3}^2 = 0$. Solving for $x_3$ (We can divide by $x_{3}^2$ because we know $x_3 \neq 0$), we get $x_3 \approx \frac{4029}{\sqrt{2014}} \approx 2\sqrt{2014} \approx 90$. We can intuitively bound $x_3$ between $88$ and $92$.

 
We will now estimate $x_1$ and $x_2$. $x_1$ and $x_2$ are close to $0$. As a result, the $\sqrt{2014}x^3$ term is negligible. We simplify and get $p(x) \approx -4029x^2 + 2 = 0$. Solving for $x$, we get $x \approx \pm \sqrt{\frac{2}{4029}} \approx \pm \sqrt{\frac{1}{2014.5}} \approx \pm \frac{1}{45}$. We can intuitively bound $x_1$ between $-\frac{1}{43}$ and $-\frac{1}{47}$. Similarly, we can intuitively bound $x_2$ between $\frac{1}{47}$ and $\frac{1}{43}$.

 
We calculate that the minimum possible value of $x_2(x_1 + x_3)$ is $\frac{88-\frac{1}{43}}{47}$ and the maximum possible value is $\frac{92-\frac{1}{47}}{43}$. The only integer that falls is this range is $\boxed{002}$.

 





\textbf{Solution 8}

Let $a=x\sqrt{2014}$. We have $\sqrt{2014}x^3-4029x^2+2=ax^2-2a^2-x^2+2=x^2(a-1)-2(a^2-1)=x^2(a-1)-2(a-1)(a+1)=(a-1)(x^2-2a-2)=0$, so $a=x\sqrt{2014}=1$ or $x^2-2a-2=x^2-2\sqrt{2014}x-2=0$. As such, the solutions are at $x=\frac{1}{\sqrt{2014}}$ and $x=\frac{2\sqrt{2014}\pm\sqrt{4\cdot2014+8}}{2}$. Note that $\frac{2\sqrt{2014}-\sqrt{4\cdot2014+8}}{2}<0<\frac{1}{\sqrt{2014}}<\sqrt{2014}<\frac{2\sqrt{2014}+\sqrt{4\cdot2014+8}}{2}$, so $x_2(x_1+x_3)=\frac{x_1+x_3}{\sqrt{2014}}=\frac{2\sqrt{2014}}{\sqrt{2014}}=\boxed{2}$. ($x_1+x_3=2\sqrt{2014}$ by vietas.)

 -zhoujef000

 





\textbf{Solution 9 (cheese) ((also wrong))}

Rewrite the polynomial as $nx^3 - (2n^2+1)x^2 + 2 = 0$. The answer probably holds for general $n$, so letting $n=0$, we have $x^2+2=0$. We'll assume that the missing root doesn't impact the final answer, so by Vietas, we have $\boxed{002}$.

 You can build confidence by setting $n=1$ to get $x^3 - 3x^2 +2 = 0 = (x-1) (x^2-2x-2)$, since there is an obvious root $x=1$, then factor the cubic.

 Middle root being the integer is only way to get an integer answer, and by Vieta, the product of the quadratic tots is the answer. You can also solve the quadratic to confirm:The roots of the quadratic are $1 \pm \sqrt{1^2 - (-2)} = 1 \pm \sqrt 3$,  so the middle root of the cubic is $1$ and the sum of the other roots is $2$. 

 You can build more confidence by setting $n=-1$ to get $-x^3 - 3x^2 +2 = 0 = (x+1) (x^2-2x^2-2)$, since there is an obvious root $x=-1$, and continue as in the $n=1$ case.

 (In fact, as other solutions show, $x=1/n$ is a root of the cubic.)

 


 
Don't use this solution. If you substitute $n=0$ into $nx^3 - (2n^2+1)x^2+2=0,$ it actually becomes $-x^2+2=0,$ not $x^2+2=0.$ So this solution is just wrong, because their answer comes out to be $-2,$ which is wrong.

 




	\item (\textit{AIME I 2014}) A disk with radius $1$ is externally tangent to a disk with radius $5$. Let $A$ be the point where the disks are tangent, $C$ be the center of the smaller disk, and $E$ be the center of the larger disk. While the larger disk remains fixed, the smaller disk is allowed to roll along the outside of the larger disk until the smaller disk has turned through an angle of $360^\circ$. That is, if the center of the smaller disk has moved to the point $D$, and the point on the smaller disk that began at $A$ has now moved to point $B$, then $\overline{AC}$ is parallel to $\overline{BD}$. Then $\sin^2(\angle BEA)=\frac{m}{n}$, where $m$ and $n$ are relatively prime positive integers. Find $m+n$.


 
 \textbf{Answer: }058



\textbf{Solution 1}


                    \begin{figure}[H]
                        \centering
                        \includegraphics[width=0.5\textwidth]{Solution/AIME_I_2014_10_Solution_Figure_1.png}
                    \end{figure}
                    

 Let $F$ be the new tangency point of the two disks. The smaller disk rolled along minor arc $\overset{\frown}{AF}$ on the larger disk.

 Let $\alpha = \angle AEF$, in radians. The smaller disk must then have rolled along an arc of length $5\alpha$, since the larger disk has a radius of $5$.  Since all of the points on major arc $\overset{\frown}{BF}$ on the smaller disk have come into contact with the larger disk at some point during the rolling, and none of the other points on the smaller disk did,
 \[\overset{\frown}{BF}=\overset{\frown}{AF}=5\alpha\]. 

 Since $\overline{AC} || \overline{BD}$, 
 \[\angle BDF \cong \angle FEA\] so the angles of minor arc $\overset{\frown}{BF}$ and minor arc $\overset{\frown}{AF}$ are equal, so minor arc $\overset{\frown}{BF}$ has an angle of $\alpha$. Since the smaller disk has a radius of $1$, the length of minor arc $\overset{\frown}{BF}$ is $\alpha$. This means that $5\alpha + \alpha$ equals the circumference of the smaller disk, so $6\alpha = 2\pi$, or $\alpha = \frac{\pi}{3}$.

 Now, to find $\sin^2{\angle BEA}$, we construct $\triangle BDE$. Also, drop a perpendicular from $D$ to $\overline{EA}$, and call this point $X$. Since $\alpha = \frac{\pi}{3}$ and $\angle DXE$ is right, 
 \[DE = 6\] 
 \[EX = 3\] and 
 \[DX = 3\sqrt{3}\] 

 Now drop a perpendicular from $B$ to $\overline{EA}$, and call this point $Y$. Since $\overline{BD} || \overline{EA}$, 
 \[XY = BD = 1\] and 
 \[BY = DX = 3\sqrt{3}\] Thus, we know that 
 \[EY = EX - XY = 3 - 1 = 2\] and by using the Pythagorean Theorem on $\triangle BEY$, we get that 
 \[BE = \sqrt{31}\] Thus, 
 \[\sin{\angle BEA} = \frac{\sqrt{27}}{\sqrt{31}}\] so 
 \[\sin^2{\angle BEA} = \frac{27}{31}\] and our answer is $27 + 31 = \boxed{058}$.

 





\textbf{Solution 2}

First, we determine how far the small circle goes. For the small circle to rotate completely around the circumference, it must rotate $5$ times (the circumference of the small circle is $2\pi$ while the larger one has a circumference of $10\pi$) plus the extra rotation the circle gets for rotating around the circle, for a total of $6$ times. Therefore, one rotation will bring point $D$ $60^\circ$ from $C$.

 Now, draw $\triangle DBE$, and call 
 \[\angle BED = x^{\circ}\] We know that $\overline{ED}$ is 6, and $\overline{BD}$ is 1. Since $EC || BD$, 
 \[\angle BDE = 60^\circ\] 

 By the Law of Cosines, 
 \[\overline{BE}^2=36+1-2\times 6\times 1\times \cos{60^\circ} = 36+1-6=31\] and since lengths are positive, 
 \[\overline{BE}=\sqrt{31}\] 

 By the Law of Sines, we know that 
 \[\frac{1}{\sin{x}}=\frac{\sqrt{31}}{\sin{60^\circ}}\] so 
 \[\sin{x} = \frac{\sin{60^\circ}}{\sqrt{31}} = \frac{\sqrt{93}}{62}\] As $x$ is clearly between $0$ and $90^\circ$, $\cos{x}$ is positive. As $\cos{x}=\sqrt{1-\sin^2{x}}$, 
 \[\cos{x} = \frac{11\sqrt{31}}{62}\] 

 Now we use the angle sum formula to find the sine of $\angle BEA$: 
 \[\sin 60^\circ\cos x + \cos 60^\circ\sin x = \frac{\sqrt{3}}{2}\frac{11\sqrt{31}}{62}+\frac{1}{2}\frac{\sqrt{93}}{62}\] 
 \[= \frac{11\sqrt{93}+\sqrt{93}}{124} = \frac{12\sqrt{93}}{124} = \frac{3\sqrt{93}}{31} = \frac{3\sqrt{31}\sqrt{3}}{31} = \frac{3\sqrt{3}}{\sqrt{31}}\]

 Finally, we square this to get 
 \[\frac{9\times 3}{31}=\frac{27}{31}\] so our answer is $27+31=\boxed{058}$.

 





\textbf{Solution 3}

We could use coordbash to solve this problem.Assume that $E$ is the origin, and $E, A, C$ are collinear points on the x-axis. Then it's trivial to see that $D=(6\cos(60^\circ),6\sin(60^\circ))$, and since $BD$ is parallel to $AC$, $B=(6\cos(60^\circ)-1,6\sin(60^\circ))$. After simplifying, we get $B=(2,3\sqrt{3})$.

 Assume there is a point $F$ such that $F=(0, 3\sqrt{3})$. Then $BF$ is parallel to $EA$, meaning $\angle FBE=\angle BEA$. $\sin(\angle FBE)$ is just $\frac{BF}{BE}=\frac{3\sqrt{3}}{\sqrt{31}}$, and squaring the fraction gives us the answer $\frac{27}{31}=\boxed{058}$.

 -Marchk26

 




	\item (\textit{AIME I 2014}) A token starts at the point $(0,0)$ of an $xy$-coordinate grid and then makes a sequence of six moves. Each move is 1 unit in a direction parallel to one of the coordinate axes. Each move is selected randomly from the four possible directions and independently of the other moves. The probability the token ends at a point on the graph of $|y|=|x|$ is $\frac{m}{n}$, where $m$ and $n$ are relatively prime positive integers. Find $m+n$.


 
 \textbf{Answer: }391



\textbf{Solution 1 (Slick Solution)}

Perform the coordinate transformation $(x, y)\rightarrow (x+y, x-y)$. Then we can see that a movement up, right, left, or down in the old coordinates adds the vectors $\langle 1, -1 \rangle$, $\langle 1, 1 \rangle$, $\langle -1, -1 \rangle$, $\langle -1, 1 \rangle$ respectively. Moreover, the transformation takes the equation $|y| = |x|$ to the union of the x and y axis. Exactly half of the moves go up in the new coordinates, and half of them go down. In order to end up on the x axis, we need to go up thrice and down thrice. The number of possible sequences of up and down moves is the number of permutations of $UUUDDD$, which is just $\binom63 = 20$. The probability of any of these sequences happening is $\left(\frac12\right)^6$. Thus, the probability of ending on the x axis is $\frac{20}{2^6}$. Similarly, the probability of ending on the y axis is the same. 

 However, we overcount exactly one case: ending at $(0, 0)$. Since ending on the x axis and ending on the y axis are independent events, the probability of both is simply $\left(\frac{20}{2^6}\right)^2 = \frac{25}{256}$. Using PIE, the total probability is $\frac{20}{64} + \frac{20}{64} - \frac{25}{256} = \frac{135}{256}$, giving an answer of $\boxed{391}$.

 


 





\textbf{Solution 2 (Casework)}

We have 4 possible moves: U, D, R, and L. The total number of paths that could be taken is $4^6$, or $4096$. There are 4 possible cases that land along the line $y = x$: $x,y = \pm 1; x,y = \pm 2; x,y = \pm 3;$ or $x = y = 0$. We will count the number of ways to end up at $(1,1), (2,2),$ and $(3,3)$, multiply them by 4 to account for the other quadrants, and add this to the number of ways to end up at $(0,0)$.

 
                \begin{itemize}
                    \item Case 1: The token ends at $(3, 3)$. In order for the token to end up here, it must have had 3 right moves, and 3 up moves. In other words, the total number of ways to get here is the ways to rearrange the letters in the sequence $RRRUUU$, which is ${6\choose 3} = 20.$
                \end{itemize}


                
                \begin{itemize}
                    \item Case 2: The token ends at $(2,2)$. In order for the token to end up here, it could have had 2 up moves, 3 right moves, and 1 left move; or 2 right moves, 3 up moves, and 1 down move. Thus, the total number of ways to get here is sum of the ways to rearrange the letters in the sequences $RRRLUU$ and $UUUDRR$, both of which are ${6\choose 1}{5\choose 2} = 60$, for a total of $120$ possibilities.
                \end{itemize}


                
                \begin{itemize}
                    <li>Case 3: The token ends at $(1,1)$. In order for the token to end up here, it could have had:
                \begin{itemize}
                    \item 1 right move, 3 up moves, and 2 down moves. In this case, the number of ways to rearrange the letters in the sequence $RUUUDD$ is ${6\choose 1}{5\choose 2} = 60.$\item 1 up move, 3 right moves, and 2 left moves. In this case, the number of ways to rearrange the letters in the sequence $URRRLL$ is ${6\choose 1}{5\choose 2} = 60.$\item 2 right moves, 1 left move, 2 up moves, and 1 down move. In this case, the number of ways to rearrange the letters in the sequence $UUDRRL$ is ${6\choose 1}{5\choose 1}{4\choose 2} = 180.$
                \end{itemize}


                </li>
                \end{itemize}


                Thus, the total number of ways to end up at $(1,1)$ is $300$.

 
                \begin{itemize}
                    <li>Case 4: The token ends at $(0,0)$. In order for the token to end up here, it could have had:
                \begin{itemize}
                    \item 3 right moves and 3 left moves. In this case, the number of ways to rearrange the letters in the sequence $RRRLLL$ is ${6\choose 3} = 20.$\item 3 up moves and 3 down moves. In this case, the number of ways to rearrange the letters in the sequence $UUUDDD$ is ${6\choose 3} = 20.$\item 1 right move, 1 left move, 2 up moves, and 2 down moves. In this case, the number of ways to rearrange the letters in the sequence $RLUUDD$ is ${6\choose 1}{5\choose 1}{4\choose 2} = 180.$\item 1 up move, 1 down move, 2 right moves, and 2 left moves. In this case, the number of ways to rearrange the letters in the sequence $RRLLUD$ is ${6\choose 1}{5\choose 1}{4\choose 2} = 180.$
                \end{itemize}


                </li>
                \end{itemize}


                Thus, the total number of ways to end up at $(0,0)$ is $400$.

 Adding these cases together, we get that the total number of ways to end up on $y = x$ is $4(20 + 120 + 300) + 400 = 2160$. Thus, our probability is $\frac{2160}{4096}$. When this fraction is fully reduced, it is $\frac{135}{256}$, so our answer is $135 + 256 = \boxed{391}.$

 





\textbf{Solution 3 (Casework)}

We split this into cases by making a chart. In each row, the entries $(\pm1)$ before the dividing line represent the right or left steps (without regard to order), and the entries after the dividing line represent the up or down steps (again, without regard to order). This table only represents the cases where the ending position $(x,y)$ satisfies $x=y$.
 \[\begin{array}{ccccccccccccl} \multicolumn{5}{c}{R (+)\qquad L (-)}& |&\multicolumn{5}{c}{U (+)\qquad D (-)}\\ +1&&  +1&&  +1&| &  +1&&  +1&&  +1\\ +1&& +1&& -1& | & +1&& +1&& -1\\ +1&& -1&&  -1& | & +1&& -1&& -1\\ -1 && -1&&  -1& |  & -1&&  -1&& -1\\ \\ +1&& +1&& +1&& -1 &|& +1&& +1\\ +1&& +1&& -1 && -1 &|& +1 && -1\\ +1&& -1&&  -1 && -1 &|& -1 &&  -1 &&(\times 2 \text{ for symmetry by swapping }R-L\text{ and }U-D)\\ \\ +1&& +1 &&+1 &&-1&& -1& |& +1\\ +1&& +1 &&-1&&  -1&& -1 &|& -1&& (\times 2\text{ symmetry})\\ \\ +1&& +1 &&+1&& -1&& -1 &&-1&| & (\times2 \text{ symmetry})\\ \end{array}\]Note that to account for the cases when $x=-y$, we can simply multiply the $U-D$ steps by $-1$, so for example, the first row would become  
 \[+1 \qquad+1\qquad +1 \ \ \ \ |\ \ \  -1\qquad -1\qquad -1.\] Therefore, we must multiply the number of possibilities in each case by $2$, except for when $x=y=0$.

 Now, we compute the number of possibilities for each case. In particular, we must compute the number of $RLUD$ words, where $R$ represents $+1$ to the left of $|$, $L$ represents $-1$ to the left of $|$, $U$ represents $+1$ to the right of $|$, and $D$ represents $-1$ to the right of $|$. Using multinomials, we compute the following numbers of possibilities for each case.
 \[{6\choose 3}\cdot  2+ \frac{6!}{2!2!}\cdot 2 + \frac{6!}{2!2!} \cdot 2 + {6\choose 3} \cdot 2 = 2(20 + 180 + 180 + 20) = 800\]
 \[\frac{6!}{3!2!}\cdot 2 + \frac{6!}{2!2!} + \frac{6!}{3!2!}\cdot 2 = 120 + 180 + 120 = 420\  (\times2\text{ for symmetry})\]
 \[\frac{6!}{3!2!} \cdot 2 + \frac{6!}{3!2!} \cdot 2 = 120 + 120 = 240\  (\times2\text{ for symmetry})\]
 \[{6\choose 3} = 20\ (\times 2\text{ for symmetry})\]

 Thus, there are $800 + 840 + 480 + 40 = 2160$ possibilities where $|x|=|y|$. Because there are $4^6$ total possibilities, the probability is $\frac{2160}{4^6} = \frac{135}{256}$, so the answer is $\boxed{391}.$

 





\textbf{Solution 4 (States)}

Denote $(x, y)_n$ the probability that starting from point $(x, y)$, the token reaches a point on the graph of $|y| = |x|$ in exactly $n$ moves. The problem asks us to find $(0, 0)_6$. We start by breaking this down:
 \[(0, 0)_6 = \frac14 \cdot ((0, 1)_5 + (0, -1)_5 + (1, 0)_5 + (-1, 0)_5)\]Notice that by symmetry, $(0, 1)_5 = (0, -1)_5 = (1, 0)_5 = (-1, 0)_5$, so the equation simplifies to
 \[(0, 0)_6 = (0, 1)_5\]We now expand $(0, 1)_5$.
 \[(0, 1)_5 = \frac14 \cdot ((0, 0)_4 + (0, 2)_4 + 2(1, 1)_4)\]First, we find $(0, 0)_4$.
 \[(0, 0)_4 = (0, 1)_3\]
 \[(0, 1)_3 = \frac14 \cdot ((0, 0)_2 + (0, 2)_2 + 2(1, 1)_2)\]At this point, we can just count the possibilities to find $(0, 0)_2 = \frac34$, $(0, 2)_2 = \frac{7}{16}$, and $(1, 1)_2 = \frac58$. Therefore,
 \[(0, 1)_3 = \frac14 \cdot (\frac34 + \frac{7}{16} + 2 \cdot \frac58)\]
 \[(0, 1)_3 = \frac{39}{64}\]Next, we find $(0, 2)_4$.
 \[(0, 2)_4 = \frac14 \cdot ((0, 1)_3 + (0, 3)_3 + 2(1, 2)_3)\]We already calculated $(0, 1)_3$, so we just need to find $(0, 3)_3$ and $(1, 2)_3$
 \[(0, 3)_3 = \frac14 \cdot ((0, 2)_2 + (0, 4)_2 + 2(1, 3)_2)\]
 \[(0, 3)_3 = \frac14 \cdot (\frac{7}{16} + 0 + 2 \cdot \frac{1}{4})\]
 \[(0, 3)_3 = \frac{15}{64}\]
 \[(1, 2)_3 = \frac14 \cdot ((1, 3)_2 + (1, 1)_2 + (0, 2)_2 + (2, 2)_2)\]
 \[(1, 2)_3 = \frac14 \cdot (\frac14 + \frac58 + \frac{7}{16} + \frac12)\]
 \[(1, 2)_3 = \frac{29}{64}\]Therefore,
 \[(0, 2)_4 = \frac14 \cdot (\frac{39}{64} + \frac{15}{64} + 2 \cdot \frac{29}{64})\]
 \[(0, 2)_4 = \frac{7}{16}\]Finally, we find $(1, 1)_4$.
 \[(1, 1)_4 = \frac12 \cdot ((0, 1)_3 + (1, 2)_3)\]
 \[(1, 1)_4 = \frac12 \cdot (\frac{39}{64} + \frac{29}{64})\]
 \[(1, 1)_4 = \frac{17}{32}\]Putting it all together,
 \[(0, 0)_6 = (0, 1)_5 =\frac14 \cdot (\frac{39}{64} + \frac{7}{16} + 2 \cdot \frac{17}{32})\]
 \[(0, 0)_6 = \frac{135}{256}\]Thus, the answer is $135 + 256 = \boxed{391}$.

 





\textbf{Solution 5 (Generating Functions)}

For any point $(x,y)$ on the coordinate plane, we can represent it as the term $a^xb^y$. Then, moving right is equivalent to multiplying by $a$, moving left is equivalent to multiplying by $1/a$, and similarly for up/down with $b$ and $1/b$. Because, on every move, the possibilities are these four directions, and each is equally likely, each move can be represented by multiplying by $a+\frac{1}{a}+b+\frac{1}{b}$. Thus, the possibilities after $6$ moves are the same as the coefficients of $\left(a+\frac{1}{a}+b+\frac{1}{b}\right)^6$. By symmetry, we just need to find the number of ways for the token to reach $(0,0),(1,1),(2,2),\text{ and }(3,3)$. By the multinomial theorem, the coefficient of $a^3b^3$ is $\binom{6}{3,3}=20$. Similarly, the coefficient of $a^2b^2$ is $2\binom{6}{3,2,1}=120$. The coefficient of $ab$ is $\binom{6}{2,2,1,1}+2\binom{6}{3,2,1}=300$. Finally, the constant term is $2\binom{6}{2,2,1,1}+2\binom{6}{3,3}=400$. Because the total number of possibilities is $4^6$, the probability is 
 \[\frac{4(20+120+300)+400}{4^6}=\frac{5+30+75+25}{4^4}=\frac{135}{256},\]and the answer is $391$.--brainiacmaniac31

 


 




	\item (\textit{AIME I 2014}) Let $A=\{1,2,3,4\}$, and $f$ and $g$ be randomly chosen (not necessarily distinct) functions from $A$ to $A$. The probability that the range of $f$ and the range of $g$ are disjoint is $\frac{m}{n}$, where $m$ and $n$ are relatively prime positive integers. Find $m$.


 
 \textbf{Answer: }453



\textbf{Solution 1 (casework)}

The natural way to go is casework. And the natural process is to sort $f$ and $g$ based on range size! Via Pigeonhole Principle, we see that the only real possibilities are: $f 1 g 1; f 1 g 2; f 1 g 3; f 2 g 2$. Note that the $1, 2$ and $1, 3$ cases are symmetrical and we need just a $*2$. Note also that the total number of cases is $4^4*4^4=4^8$.

 $f 1 g 1$: clearly, we choose one number as the range for $f$, one for $g$, yielding $12$ possibilities.

 $f 1 g 2$ with symmetry (WLOG $f$ has 1 element): start by selecting numbers for the ranges. This yields $4$ for the one number in $f$, and $3$ options for the two numbers for $g$. Afterwards, note that the function with 2 numbers in the range can have $4+6+4=14$ arrangements of these two numbers (1 of one, 3 of the other *2 and 2 of each). Therefore, we have $2*12*14$ possibilities, the 2 from symmetry.

 $f 2 g 2$: no symmetry, still easy! Just note that we have $6$ choices of which numbers go to $f$ and $g$, and within each, $14*14=196$ choices for the orientation of each of the two numbers. That's $6*196$ possibilities.

 $f 1 g 3$: again, symmetrical (WLOG $f$ has one element): $4$ ways to select the single element for $f$, and then find the number of ways to distribute the $3$ distinct numbers in the range for $g$. The only arrangement for the frequency of each number is ${1, 1, 2}$ in some order. Therefore, we have $3$ ways to choose which number is the one represented twice, and then note that there are $12$ ways to arrange these! The number of possibilities in this situation is $2 * 4 * 3 * 12$.

 Total, divided by $4^8$, gets $\frac{3 * (1 + 2 * 7^2 + 2^2 * 7 + 2^3 * 3)}{4^7}$, with numerator $\boxed{453}$.

 





\textbf{Solution 2 (casework)}

We note there are $4^4 = 256$ possibilities for each of $f$ and $g$ from $A$ to $A$ since the input of the four values of each function has four options each for an output value.

 We proceed with casework to determine the number of possible $f$ with range 1, 2, etc.

 
                \begin{itemize}
                    \item Range 1:
                \end{itemize}


                There are 4 possibilities: all elements output to 1, 2, 3, or 4.

 
                \begin{itemize}
                    \item Range 2:
                \end{itemize}


                We have ${{4}\choose {2}} = 6$ ways to choose the two output elements for $f$. At this point we have two possibilities: either $f$ has 3 of 1 element and 1 of the other, or 2 of each element. In the first case, there are 2 ways to pick the element which there are 3 copies of, and ${{4}\choose {1}} = 4$ ways to rearrange the 4 elements, for a total of $6 * 4 * 2 = 48$ ways for this case. For the second case, there are ${{4}\choose {2}} = 6$ ways to rearrange the 4 elements, for a total of $6 * 6 = 36$ ways for this case. Adding these two, we get a total of $36 + 48 = 84$ total possibilities.

 
                \begin{itemize}
                    \item Range 3:
                \end{itemize}


                We have ${{4}\choose {3}} = 4$ ways to choose the three output elements for $f$. We know we must have 2 of 1 element and 1 of each of the others, so there are 3 ways to pick this element. Finally, there are ${{4}\choose{1}}*{{3}\choose{1}} = 12$ ways to rearrange these elements (since we can pick the locations of the 2 single elements in this many ways), and our total is $4 * 3 * 12 = 144$ ways.

 
                \begin{itemize}
                    \item Range 4:
                \end{itemize}


                Since we know the elements present, we have $4!$ ways to arrange them, or 24 ways.

 (To check, $4 + 84 + 144 + 24 = 256$, which is the total number of possibilities).

 We now break $f$ down by cases, and count the number of $g$ whose ranges are disjoint from $f$'s.

 
                \begin{itemize}
                    \item Case 1: $f$'s range contains 1 element
                \end{itemize}


                We know that there are 3 possibilities for $g$ with 1 element. Since half the possibilities for $g$ with two elements will contain the element in $f$, there are $84/2 = 42$ possibilities for $g$ with 2 elements. Since $3/4$ the possibilities for $g$ with 3 elements will contain the element in $f$, there are $144/4 = 36$ possibilities for $g$ with 3 elements. Clearly, no 4-element range for $g$ is possible, so the total number of ways for this case to happen is $4(3 + 42 + 36) = 324$.

 
                \begin{itemize}
                    \item Case 2: $f$'s range contains 2 elements
                \end{itemize}


                We know that there are 2 possibilities for $g$ with 1 element. If $g$ has 2 elements in its range, they are uniquely determined, so the total number of sets with a range of 2 elements that work for $g$ is $84/6 = 14$. No 3-element or 4-element ranges for $g$ are possible. Thus, the total number of ways for this to happen is $84(2 + 14) = 1344$.

 


 
                \begin{itemize}
                    \item Case 3:  $f$'s range contains 3 elements
                \end{itemize}


                In this case, there is only 1 possibility for $g$ - all the output values are the element that does not appear in $f$'s range. Thus, the total number of ways for this to happen is $144$.

 
                \begin{itemize}
                    \item Summing the cases
                \end{itemize}


                We find that the probability of $f$ and $g$ having disjoint ranges is equal to:

 $\dfrac{324 + 1344 + 144}{256^2}=\dfrac{1812}{2^{16}}= \dfrac{453}{2^{14}}$

 Thus, our final answer is $\boxed{453}$.

 





\textbf{Solution 3 (simplification of Solution 2)}

As before, there are 4 functions with a range of size 1, 84 with a range of size 2, and 144 with a range of size 3. If the range of $f$ has size $k$, the codomain of $g$ is restricted to a set of size $4 - k$. Any function from $A$ into this codomain will do, so there are $(4 - k)^4$ possibilities for $g$ given a function $f$. The probability of $f$ and $g$ having disjoint ranges is then
 \[\frac{4\cdot 3^4 + 84\cdot 2^4 + 144\cdot 1^4}{(4^4)^2} = \frac{\boxed{453}}{2^{14}}.\]

 


 


 




	\item (\textit{AIME I 2014}) On square $ABCD$, points $E,F,G$, and $H$ lie on sides $\overline{AB},\overline{BC},\overline{CD},$ and $\overline{DA},$ respectively, so that $\overline{EG} \perp \overline{FH}$ and $EG=FH = 34$. Segments $\overline{EG}$ and $\overline{FH}$ intersect at a point $P$, and the areas of the quadrilaterals $AEPH, BFPE, CGPF,$ and $DHPG$ are in the ratio $269:275:405:411.$ Find the area of square $ABCD$.


 
                    \begin{figure}[H]
                        \centering
                        \includegraphics[width=0.5\textwidth]{Problem/AIME_I_2014_13_Problem_Figure_1.png}
                    \end{figure}
                    


 
 \textbf{Answer: }850



\textbf{Solution (Official Solution, MAA)}

Let $s$ be the side length of $ABCD$, let $Q$, and $R$ be the midpoints of $\overline{EG}$ and $\overline{FH}$, respectively, let $S$ be the foot of the perpendicular from $Q$ to $\overline{CD}$, let $T$ be the foot of the perpendicular from $R$ to $\overline{AD}$. 
                    \begin{figure}[H]
                        \centering
                        \includegraphics[width=0.5\textwidth]{Solution/AIME_I_2014_13_Solution_Figure_1.png}
                    \end{figure}
                    The fraction of the area of the square $ABCD$ which is occupied by trapezoid $BCGE$ is 
 \[\frac{275+405}{269+275+405+411}=\frac 12,\]so $Q$ is the center of $ABCD$. Thus $R$, $Q$, $S$ are collinear, and $RT=QS=\frac 12 s$. Similarly, the fraction of the area occupied by trapezoid $CDHF$ is $\frac 35$, so $RS=\frac 35s$ and $RQ=\frac{1}{10}s$. 

 Because $\triangle QSG \cong \triangle RTH$, the area of $DHPG$ is the sum 
 \[[DHPG]=[DTRS]+[RPQ].\] Rectangle $DTRS$ has area $RS\cdot RT = \frac 35s\cdot \frac 12 s = \frac{3}{10}s^2$. If $\angle QRP = \theta$ , then $\triangle RPQ$ has area 
 \[[RPQ]= \frac 12 \cdot \frac 1{10}s\sin\theta \cdot \frac 1{10}s\cos\theta = \frac 1{400}s^2\sin 2\theta.\]Therefore the area of $[DHPG]$ is $s^2(\frac 3{10}+\frac 1{400}\sin 2\theta)$. Because the area of trapezoid $CDHF$ is $\frac 35 s^2$, the area of $CGPF$ is $s^2(\frac 3{10}-\frac 1{400}\sin 2\theta)$.

 Because these areas are in the ratio $411:405=(408+3):(408-3)$, it follows that 
 \[\frac{\frac 1{400}\sin 2\theta}{\frac 3{10}}=\frac 3{408},\]from which we get $\sin 2\theta = \frac {15}{17}$. Note that $\theta =\angle RHT > \angle QAT = 45^\circ$, so $\cos 2\theta = -\sqrt{1-\sin^2 2\theta}= -\frac 8{17}$ and $\sin^2\theta = \frac{1}{2}(1-\cos 2\theta) = \frac{25}{34}$. Then 
 \[[ABCD]=s^2 = EG^2\sin^2\theta = 34^2 \cdot \frac {25}{34} = 850.\]

 





\textbf{Solution 1}

Let $s$ be the side length of $ABCD$, let $[ABCD]=1360a$. Let $Q$ and $R$ be the midpoints of $\overline{EG}$ and $\overline{FH}$, respectively; because $269+411=275+405$, $Q$ is also the center of the square. Draw $\overline{IJ} \parallel \overline{HF}$ through $Q$, with $I$ on $\overline{AD}$, $J$ on $\overline{BC}$. 
                    \begin{figure}[H]
                        \centering
                        \includegraphics[width=0.5\textwidth]{Solution/AIME_I_2014_13_Solution_Figure_2.png}
                    \end{figure}
                    Segments $\overline{EG}$ and $\overline{IJ}$ divide the square into four congruent quadrilaterals, each of area $\frac 14 [ABCD]=340a$. Then 
 \[[HFJI]=[ABJI]-[ABFH]=136a.\] The fraction of the total area occupied by parallelogram $HFJI$ is $\frac 1{10}$, so $RQ=\frac{s}{10}$. 

 Because $[HFJI]= HF\cdot PQ$, with $HF=34$, we get $PQ=4a$. Now 
 \[[PQR]=[HPQI]-[HRQI]= ([AEQI]-[AEPH])-\frac 12[IJFH] = 3a,\] and because $[PQR]=\frac 12 \cdot PQ\cdot PR$, with $PQ=4a$, we get $PR=\frac 32$. By Pythagoras' Theorem on $\triangle PQR$, we get 
 \[16a^2+\frac 94 =\frac{68}{5}a,\quad \text{i.e.}\quad 320a^2-272a+45=0,\] with roots $a=\frac 9{40}$ or $a=\frac58$. The former leads to a square with diagonal less than $34$, which can't be, since $EG=FH=34$; therefore $a=\frac 58$ and $[ABCD]=850$.

 





\textbf{Solution 2 (Fakesolve)}

$269+275+405+411=1360$, a multiple of $17$. In addition, $EG=FH=34$, which is $17\cdot 2$.Therefore, we suspect the square of the "hypotenuse" of a right triangle, corresponding to $EG$ and $FH$  must be a multiple of $17$. All of these triples are primitive:

 
 \[17=1^2+4^2\]
 \[34=3^2+5^2\]
 \[51=\emptyset\]
 \[68=\emptyset\text{ others}\]
 \[85=2^2+9^2=6^2+7^2\]
 \[102=\emptyset\]
 \[119=\emptyset \dots\]

 The sides of the square can only equal the longer leg, or else the lines would have to extend outside of the square. Substituting $EG=FH=34$:
 \[\sqrt{17}\rightarrow 34\implies 8\sqrt{17}\implies A=\textcolor{red}{1088}\]
 \[\sqrt{34}\rightarrow 34\implies 5\sqrt{34}\implies A=850\]
 \[\sqrt{85}\rightarrow 34\implies \{18\sqrt{85}/5,14\sqrt{85}/5\}\implies A=\textcolor{red}{1101.6,666.4}\]

 Thus, $\boxed{850}$ is the only valid answer.

 





\textbf{Solution 3}

Continue in the same way as solution 1 to get that $POK$ has area $3a$, and $OK = \frac{d}{10}$. You can then find $PK$ has length $\frac 32$.

 Then, if we drop a perpendicular from $H$ to $BC$ at $L$, We get $\triangle HLF \sim \triangle OPK$.

 Thus, $LF = \frac{15\cdot 34}{d}$, and we know $HL = d$, and $HF = 34$. Thus, we can set up an equation in terms of $d$ using the Pythagorean theorem.

 
 \[\frac{15^2 \cdot 34^2}{d^2} + d^2 = 34^2\]

 
 \[d^4 - 34^2 d^2 + 15^2 \cdot 34^2 = 0\]

 
 \[(d^2 - 34 \cdot 25)(d^2 - 34 \cdot 9) = 0\]

 $d^2 = 34 \cdot 9$ is extraneous, so $d^2 = 34 \cdot 25$. Since the area is $d^2$, we have it is equal to $34 \cdot 25 = \boxed{850}$

 -Alexlikemath

 





\textbf{Solution 4}

saw this at bottom of \href{https://artofproblemsolving.com/community/c5h580667p3428922}{https://artofproblemsolving.com/community/c5h580667p3428922}

 
                    \begin{figure}[H]
                        \centering
                        \includegraphics[width=0.5\textwidth]{Solution/AIME_I_2014_13_Solution_Figure_3.png}
                    \end{figure}
                    

 
Let $O$ be the center of the square. Then let $Q$ be on $HF$ susch that $OQ \perp CD$. Draw $H'F' \parallel HF$ through $O$. Let $s$ be the side length of the square, so that $s^2$ is the desired.

 First, we note that as $w + z = x + y$, the area of $[ADGC]$ is half the area of $[ABCD]$, which means $O \in EG$. Similarly, $[CDHF] = y + z = \frac{3}{5}[ABCD]$ because of the ratios. But because $O \in H'F'$, we have $[CDH'F'] = \frac{1}{2}[ABCD]$. Subtracting yields $[H'HFF'] = \frac{1}{10}[ABCD]$. But as it is a parallelogram with base $34$ and height $OP$, we get:
 \[340(OP) = s^2\]Now consider quadrilateral $OEBF'$, which has area $1/4$ of the square. We can write:
 \[\frac{1}{4}s^2 = [OEBF'] = [PEBF] + [FPOF'] = x + [FQOF'] - [QOP]\]However, we know the areas of $x$ and parallelogram $[FQOF'] = 17OP = \frac{s^2}{20}$ by our earlier result. Substituting these yields:

 
 \[[QOP] = \frac{3}{1360}s^2\]But as we know $OP = \frac{s^2}{340}$, we get $QP = 3/2$.Because $CDHF$ is a trapezoid of height $s$ and area $\frac{3}{5}s^2$, we have $\frac{1}{2}s(CF + DH) = \frac{3}{5}s^2$. This means:
 \[CF + DH = \frac{6}{5}s\]But using the same method on trapezoid $CDH'F'$ of area $\frac{1}{2}s^2$, we get
 \[CF' + DH' = s\]And subtracting yields $HH' + FF' = 2OQ = \frac{1}{5}s$, which gives $OQ = \frac{s}{10}$. At this point, applying Pythagoras on $\triangle QOP$ yields:
 \[\left( \frac{s^2}{340} \right)^2 + \left( \frac{3}{2} \right)^2 = \left( \frac{s}{10} \right)^2\]and the resulting quadratic is easy to solve for $s^2 = 306, 850$. But as $[ABCD] > [EFGH] = 34^2/2 > 306$, we have $\boxed{850}$ as the solution.

 





\textbf{Solution 5 (Coordinate Geometry)}

Let $s$ be the side length of $ABCD$. We see that quadrilaterals $BCGE$ and $ADGE$ each occupies half the area of $ABCD$. This follows that $BE=DG$. Let $a=AF,b=BE,c=AH$. We establish a coordiante system with $B$ being the origin. Here is our setup:\begin{align*}A&=(0,s)\\B&=(0,0)\\C&=(s,0)\\D&=(s,s)\\E&=(0,b)\\F&=(a,0)\\G&=(s,s-b)\\H&=(c,s)\\P&=(p,q)\\FH&:y=-\frac{s}{a-c}(x-a)\end{align*}
Let $\omega,x,y,z$ be the areas of the quadrilaterals $AEPH, BFPE, CGPF,$ and $DHPG$. By the Shoelace formula, we have\begin{align*}\omega&=\frac{1}{2}\Bigl(c(s-q)+p(s-b)\Bigr)\\x&=\frac{1}{2}(aq+bp)\\y&=\frac{1}{2}\Bigl((s-p)(s-b)+q(s-a)\Bigr)\\z&=\frac{1}{2}\Bigl((s-c)(s-q)+b(s-p)\Bigr)\end{align*}
Let $H'$ be the feet of perpendicular from $H$ to $\overline{BC}$, and we define $G'$ similarly. From $\triangle HH'F \cong \triangle GG'E$, we have $H'F=G'E\implies a-c=s-2b$, and let this distance be $k$. To simplify the system of equations, we make the following substitutions:
 \[u=\frac{k}{s},\qquad p'=\frac{p}{s},\qquad q'=\frac{q}{s}\]By the Pythagorean's theorem, $s^2+k^2=34^2\implies s^2=\frac{34^2}{1+u^2}$. We need to find the value of $u$.

 The equation of $FH$ becomes $y=-\frac{1}{u}(x-a)$. Substituting the coordinates of $P=(p,q)$, we have
 \[a=p+uq\]Therefore,\begin{align*}\frac{b}{s}&=\frac{1-u}{2}\\\frac{a}{s}&=p'+u\cdot q'\\\frac{c}{s}&=p'+u\cdot\left(q'-1\right)\\\end{align*}
Those expressions are especially helpful when we solve the following equation:
 \[y-\omega=\frac{405-269}{269+275+405+411}s^2=\frac{1}{10}s^2\]We proceed by computing the above equation:\begin{align*}y-\omega&=\frac{1}{2}\Bigl((s-2p)(s-b)+q(s-a)-c(s-q)\Bigr)\\\frac{y-\omega}{s^2}&=\frac{1}{2}\left(\left(1-\frac{2p}{s}\right)\left(1-\frac{b}{s}\right)+\frac{q}{s}-\frac{q}{s}\left(\frac{a}{s}-\frac{c}{s}\right)-\frac{c}{s}\right)\\\frac{1}{10}&=\frac{1}{2}\left((1-2p')\cdot\frac{1+u}{2}+q'-uq'-p'-uq'+u\right)\\\frac{2}{5}&=(1-2p')(1+u)+2q'-2uq'-2p'-2uq'+2u\\\frac{2}{5}&=1+3u-(4+2u)p'+(2-4u)q'\tag{1}\end{align*}
We compute $\frac{z-x}{s^2}=\frac{1}{10}$ similarly:\begin{align*}z-x&=\frac{1}{2}\Bigl(b(s-2p)+(s-c)(s-q)-aq\Bigr)\\\frac{z-x}{s^2}&=\frac{1}{2}\left(\frac{b}{s}\left(1-\frac{2p}{s}\right)+1-\frac{c}{s}-\frac{q}{s}-\frac{q}{s}\left(\frac{a}{s}-\frac{c}{s}\right)\right)\\\frac{1}{10}&=\frac{1}{2}\left(\frac{1-u}{2}\cdot(1-2p')+1-p'-uq'+u-q'-uq'\right)\\\frac{2}{5}&=(1-u)(1-2p')+2-2p'-2uq'+2u-2q'-2uq'\\\frac{2}{5}&=3+u-(4-2u)p'-(2+4u)q'\tag{2}\end{align*}
Rewriting $(1)$ and $(2)$, we have\begin{align*}10(2+u)p'-10(1-2u)q'&=3(1+5u)\\10(2-u)p'+10(1+2u)q'&=13+5u\end{align*}
We solve $p',q'$ in terms of $u$ by the Cramer's Rule:
 \[p'=\frac{4+5u^2}{10(1+u^2)},\qquad q'=\frac{5-u+5u^2}{10(1+u^2)}\]We proceed by solving $\frac{x}{s^2}=\frac{55}{272}$:\begin{align*}x&=\frac{1}{2}(aq+bp)\\\frac{x}{s^2}&=\frac{1}{2}\left(\frac{a}{s}\cdot\frac{q}{s}+\frac{b}{s}\cdot\frac{p}{s}\right)\\\frac{55}{272}&=\frac{1}{2}\Bigl((p'+uq')q'+\frac{1-u}{2}\cdot p'\Bigr)\\\frac{55}{136}&=(p'+uq')q'+\frac{1-u}{2}\cdot p'\end{align*}
Pluging in the expressions for $p',q'$, we have
 \[\frac{55}{136}=\frac{4+5u+4u^2+5u^3}{10(1+u^2)}\cdot\frac{5-u+5u^2}{10(1+u^2)}+\frac{4-4u+5u^2-5u^3}{20(1+u^2)}\]Or
 \[\frac{40+u+80u^2+u^3+40u^4}{25(1+2u^2+u^4)}=\frac{55}{34}\]Recall that $u=\frac{k}{s}\in (0,1)$. Solving the palindrome equation gives $u=\frac{3}{5}$. Therefore,
 \[s^2=\frac{34^2}{1+u^2}=34^2\cdot\frac{25}{34}=\boxed{850}\]

 


 




	\item (\textit{AIME I 2014}) Let $m$ be the largest real solution to the equation


 
 \[\dfrac{3}{x-3} + \dfrac{5}{x-5} + \dfrac{17}{x-17} + \dfrac{19}{x-19} = x^2 - 11x - 4\]


 There are positive integers $a, b,$ and $c$ such that $m = a + \sqrt{b + \sqrt{c}}$. Find $a+b+c$.


 
 \textbf{Answer: }263



\textbf{Solution 1}

The first step is to notice that the 4 on the right hand side can simplify the terms on the left hand side. If we distribute 1 to $\frac{3}{x-3}$, then the fraction becomes of the form $\frac{x}{x - 3}$. A similar cancellation happens with the other four terms. If we assume $x = 0$ is not the highest solution (which is true given the answer format) we can cancel the common factor of $x$ from both sides of the equation.

 
 \[\frac{1}{x - 3} + \frac{1}{x - 5} + \frac{1}{x - 17} + \frac{1}{x - 19} = x - 11\]

 Then, if we make the substitution $y = x - 11$, we can further simplify.

 
 \[\frac{1}{y + 8} + \frac{1}{y + 6} + \frac{1}{y - 6} + \frac{1}{y - 8} =y\]

 If we group and combine the terms of the form  $y - n$ and   $y + n$, we get this equation:

 
 \[\frac{2y}{y^2 - 64} + \frac{2y}{y^2 - 36} = y\]

 Then, we can cancel out a $y$ from both sides, knowing that $x = 11$ is not a possible solution given the answer format.After we do that, we can make the final substitution $z = y^2$.

 
 \[\frac{2}{z - 64} + \frac{2}{z - 36} = 1\]

 
 \[2z - 128 + 2z - 72 = (z - 64)(z - 36)\]

 
 \[4z -  200 = z^2 - 100z + 64(36)\]

 
 \[z^2 - 104z + 2504 = 0\]

 Using the quadratic formula, we get that the largest solution for $z$ is $z = 52 + 10\sqrt{2}$. Then, repeatedly substituting backwards, we find that the largest value of $x$ is $11 + \sqrt{52 + \sqrt{200}}$. The answer is thus $11 + 52 + 200 = \boxed{263}$

 
Note: When $x$ is barely larger than $19$, then $\frac{19}{x-19}$ is very large, so the left side of the equation approaches infinity as $x$ approaches $19$ from the side greater than $19$. However, we also know as $x$ gets very large, the fractions get smaller as the left side approaches $0$. Since the quadratic on the right side is increasing and positive when $x=19$, the equation will be true at a certain $x>19.$ So, we don't have to assume there is an answer $x>0.$

 





\textbf{Solution 2}

Proceed as with Solution 1 until we get the following.

 $\frac{2y}{y^2 - 64} + \frac{2y}{y^2 - 36} = y \implies$

 $\frac{1}{y^2 - 64} + \frac{1}{y^2 - 36} = \frac{1}{2}.$

 Here, we may also use a slightly different substitution, $z = y^2 - 50.$ This gives:

 $\frac{1}{z - 14} + \frac{1}{z + 14} = \frac{1}{2} \implies$

 $\frac{2z}{z^2 - 196} = \frac{1}{2} \implies$

 $z^2 - 4z - 196 = 0.$

 We now have a simpler quadratic, eliminating tedious and potentially error-prone calculations. Following through, we get $z = 2 + \sqrt{200} \implies y = \sqrt{52 + \sqrt{200}}$ as desired.

 





\textbf{Solution 3}

Stolen from forums

 \begin{align}& \frac{3}{x-3}+\frac{5}{x-5}+\frac{17}{x-17}+\frac{19}{x-19}={{x}^{2}}-11x-4 \\& x-11=u\Rightarrow x=u+11 \\& \Rightarrow \frac{3}{u+8}+\frac{5}{u+6}+\frac{17}{u-6}+\frac{19}{u-8}=(u+11)u-4 \\& \Rightarrow \frac{3(u-8)+19(u+8)}{{{u}^{2}}-64}+\frac{5(u-6)+17(u+6)}{{{u}^{2}}-36}={{u}^{2}}+11u-4 \\& \Rightarrow \frac{22u+128}{{{u}^{2}}-64}+\frac{22u+72}{{{u}^{2}}-36}={{u}^{2}}+11u-4 \\& \Rightarrow \frac{2({{u}^{2}}+11u)}{{{u}^{2}}-64}+\frac{2({{u}^{2}}+11u)}{{{u}^{2}}-36}={{u}^{2}}+11u \\& \Rightarrow \left\{ \begin{matrix}{{u}^{2}}+11u=0\Rightarrow u=0,-11 \\\frac{2}{{{u}^{2}}-64}+\frac{2}{{{u}^{2}}-36}=1\Rightarrow ({{u}^{2}}-36)({{u}^{2}}-64)=4{{u}^{2}}-200 \\\end{matrix} \right. \\& {{u}^{4}}-104{{u}^{2}}+2504=0\Rightarrow {{u}^{2}}=52\pm 10\sqrt{2} \\& \Rightarrow u=\pm \sqrt{52\pm 10\sqrt{2}}\Rightarrow x=11\pm \sqrt{52\pm 10\sqrt{2}}\end{align}

 





\textbf{Solution 4}

We first consider adding $1$ to each fraction, which yields $\frac{x}{x-3} + \frac{x}{x-5} + \frac{x}{x-17} + \frac{x}{x-19} = x^2-11x$. We now divide by $x$ on both sides to yield $\frac{1}{x-3}+\frac{1}{x-5}+\frac{1}{x-17}+\frac{1}{x-19} = x-11$. Next, consider grouping the $\frac{1}{x-3}$ and $\frac{1}{x-19}$, as well as the $\frac{1}{x-5}$ and $\frac{1}{x-17}$. If we combine these, we get $\frac{2x-22}{(x-3)(x-19)} + \frac{2x-22}{(x-15)(x-7)} = x-11$. Now, let us divide by $2x-22$ on both sides to yield $\frac{1}{(x-3)(x-19)} + \frac{1}{(x-15)(x-7)} = \frac{1}{2}$. We $\frac{(x-3)(x-19)+(x-15)(x-17)}{(x-3)(x-19)(x-5)(x-17)} = \frac{1}{2}$. Doing a bit of algebra bash yields $2(2x^2-44x+142)$ on the LHS after also cross multiplying. As for the RHS, we get $(x^2-22x+85)(x^2-22x+57)$. Now, let $a=x^2-22x+71$. We get $2(2a)=(a-14)(a+14)$, or $a^2-4a-196$. Thus, $a=2+10\sqrt2$ or $a=2-10\sqrt2$. We drop the latter as it is smaller. Next, we get $2+10\sqrt{2}=x^2-22x+71$, or $x^2-22x+(69-10\sqrt2)$. We now complete the square to get $(x-11)^2 - 52 -10\sqrt2=0$. Bringing that to the RHS gives $(x-11)^2=52+10\sqrt2$, so $x=11+\sqrt{52+10\sqrt{2}}$. However, we need to expand the last square root to get $x=11+\sqrt{52+\sqrt{200}}$. Thus, $a+b+c=11+52+200=263$.

 




	\item (\textit{AIME I 2014}) In $\triangle ABC, AB = 3, BC = 4,$ and $CA = 5$. Circle $\omega$ intersects $\overline{AB}$ at $E$ and $B, \overline{BC}$ at $B$ and $D,$ and $\overline{AC}$ at $F$ and $G$. Given that $EF=DF$ and $\frac{DG}{EG} = \frac{3}{4},$ length $DE=\frac{a\sqrt{b}}{c},$ where $a$ and $c$ are relatively prime positive integers, and $b$ is a positive integer not divisible by the square of any prime. Find $a+b+c$.


 
 \textbf{Answer: }041



\textbf{Solution 1}

Since $\angle DBE = 90^\circ$, $DE$ is the diameter of $\omega$. Then $\angle DFE=\angle DGE=90^\circ$. But $DF=FE$, so $\triangle DEF$ is a 45-45-90 triangle. Letting $DG=3x$, we have that $EG=4x$, $DE=5x$, and $DF=EF=\frac{5x}{\sqrt{2}}$. 

 Note that $\triangle DGE \sim \triangle ABC$ by SAS similarity, so $\angle BAC = \angle GDE$ and $\angle ACB = \angle DEG$. Since $DEFG$ is a cyclic quadrilateral, $\angle BAC = \angle GDE=180^\circ-\angle EFG = \angle AFE$ and $\angle ACB = \angle DEG = \angle GFD$, implying that $\triangle AFE$ and $\triangle CDF$ are isosceles. As a result, $AE=CD=\frac{5x}{\sqrt{2}}$, so $BE=3-\frac{5x}{\sqrt{2}}$ and $BD =4-\frac{5x}{\sqrt{2}}$. 

 Finally, using the Pythagorean Theorem on $\triangle BDE$, 
 \[\left(3-\frac{5x}{\sqrt{2}}\right)^2 + \left(4-\frac{5x}{\sqrt{2}}\right)^2 = (5x)^2\]Solving for $x$, we get that $x=\frac{5\sqrt{2}}{14}$, so $DE= 5x =\frac{25 \sqrt{2}}{14}$. Thus, the answer is $25+2+14=\boxed{041}$.

 





\textbf{Solution 2}


                    \begin{figure}[H]
                        \centering
                        \includegraphics[width=0.5\textwidth]{Solution/AIME_I_2014_15_Solution_Figure_1.png}
                    \end{figure}
                    

 First we note that $\triangle DEF$ is an isosceles right triangle with hypotenuse $\overline{DE}$ the same as the diameter of $\omega$. We also note that $\triangle DGE \sim \triangle ABC$ since $\angle EGD$ is a right angle and the ratios of the sides are $3:4:5$. 

 From congruent arc intersections, we know that $\angle GED \cong \angle GBC$, and that from similar triangles $\angle GED$ is also congruent to $\angle GCB$. Thus, $\triangle BGC$ is an isosceles triangle with $BG = GC$, so $G$ is the midpoint of $\overline{AC}$ and $AG = GC = 5/2$. Similarly, we can find from angle chasing that $\angle ABF = \angle EDF = \frac{\pi}4$. Therefore, $\overline{BF}$ is the angle bisector of $\angle B$. From the angle bisector theorem, we have $\frac{AF}{AB} = \frac{CF}{CB}$, so $AF = 15/7$ and $CF = 20/7$. 

 Lastly, we apply power of a point from points $A$ and $C$ with respect to $\omega$ and have $AE \times AB=AF \times AG$ and $CD \times CB=CG \times CF$, so we can compute that $EB = \frac{17}{14}$ and $DB = \frac{31}{14}$. From the Pythagorean Theorem, we result in $DE = \frac{25 \sqrt{2}}{14}$, so $a+b+c=25+2+14= \boxed{041}$

 
Also: $FG=\frac{20}{7}-\frac{5}{2}=\frac{5}{2}-\frac{15}{7}=\frac{5}{14}$. We can also use Ptolemy's Theorem on quadrilateral $DEFG$ to figure what $FG$ is in terms of $d$:
 \[DE\cdot FG+DG\cdot EF=DF\cdot EG\]
 \[d\cdot FG+\frac{3d}{5}\cdot \frac{d}{\sqrt{2}}=\frac{4d}{5}\cdot \frac{d}{\sqrt{2}}\]
 \[d\cdot FG+\frac{3d^2}{5\sqrt{2}}=\frac{4d^2}{5\sqrt{2}}\implies FG=\frac{d}{5\sqrt{2}}\]Thus $\frac{d}{5\sqrt{2}}=\frac{5}{14}\rightarrow d=5\sqrt{2}\cdot\frac{5}{14}=\frac{25\sqrt{2}}{14}$. $a+b+c=25+2+14= \boxed{041}$

 





\textbf{Solution 3}

Call $DE=x$ and as a result $DF=EF=\frac{x\sqrt{2}}{2}, EG=\frac{4x}{5}, GD=\frac{3x}{5}$. Since $EFGD$ is cyclic we just need to get $DG$ and using LoS(for more detail see the $2$nd paragraph of Solution $2$) we get $AG=\frac{5}{2}$ and using a similar argument(use LoS again) and subtracting you get $FG=\frac{5}{14}$ so you can use Ptolemy to get $x=\frac{25\sqrt{2}}{14} \implies \boxed{041}$.~

 





\textbf{Solution 4}

See inside the $\triangle DEF$, we can find that $AG>AF$ since if $AG<AF$, we can see that Ptolemy Theorem inside cyclic quadrilateral $EFGD$ doesn't work. Now let's see when $AG>AF$, since  $\frac{DG}{EG} = \frac{3}{4}$, we can assume that $EG=4x;GD=3x;ED=5x$, since we know $EF=FD$ so $\triangle EFD$ is isosceles right triangle. We can denote $DF=EF=\frac{5x\sqrt{2}}{2}$.Applying Ptolemy Theorem inside the cyclic quadrilateral $EFGD$ we can get the length of $FG$ can be represented as $\frac{x\sqrt{2}}{2}$. After observing, we can see $\angle AFE=\angle EDG$, whereas $\angle A=\angle EDG$ so we can see $\triangle AEF$ is isosceles triangle. Since $\triangle ABC$ is a $3-4-5$ triangle so we can directly know that the length of AF can be written in the form of $3x\sqrt{2}$. Denoting a point $J$ on side $AC$ with that $DJ$ is perpendicular to side $AC$. Now with the same reason, we can see that $\triangle DJG$ is a isosceles right triangle, so we can get $GJ=\frac{3x\sqrt{2}}{2}$ while the segment $CJ$ is $2x\sqrt{2}$ since its 3-4-5 again. Now adding all those segments together we can find that $AC=5=7x\sqrt{2}$ and $x=\frac{5\sqrt{2}}{14}$ and the desired $ED=5x=\frac{25\sqrt{2}}{14}$which our answer is $\boxed{041}$ ~bluesoul

 





\textbf{Solution 5}


                    \begin{figure}[H]
                        \centering
                        \includegraphics[width=0.5\textwidth]{Solution/AIME_I_2014_15_Solution_Figure_2.png}
                    \end{figure}
                    The main element of the solution is the proof that $BF$ is bisector of $\angle B.$

 Let $O$ be the midpoint of $DE.$ $\angle EBF = 90^\circ \implies$

 $O$ is the center of the circle $BDGFE.$ $\angle EOF = 90^\circ \implies \overset{\Large\frown} {EF} = 90^\circ \implies \angle EBF = 45^\circ \implies$BF is bisector of $\angle ABC\implies BF = \frac {2AB \cdot BC}{AB+BC} \cos 45^\circ =\frac {12 \cdot \sqrt{2}}{7}.$ 
 \[\angle EGD = 90^\circ, \frac {EG}{GD}=\frac{4}{3} \implies\]
 \[\angle GED = \angle GCD =\gamma \implies  \overset{\Large\frown} {DG} = 2\gamma.\]
 \[2\angle ACB =  \overset{\Large\frown} {BEF} -  \overset{\Large\frown} {DG} \implies  \overset{\Large\frown} {BEF} = 4 \gamma \implies\]
 \[\angle BOF = 4 \gamma \implies \angle OBF = \angle OFB = 90^\circ – 2 \gamma.\]Let $BO = EO = DO = r \implies BF = 2 r \cos(90^\circ – 2\gamma) =$
 \[=2 r \sin 2\gamma = 4r \sin \gamma \cdot \cos \gamma = 4 r\cdot \frac {3}{5} \cdot \frac {4}{5} = \frac {48}{25} = \frac {12 \cdot \sqrt{2}}{7}\implies\]
 \[r = \frac {25 \cdot \sqrt{2}}{28}\implies ED = 2r =  \frac {25 \cdot \sqrt{2}}{14}\implies  \boxed{\textbf{041}}.\]\textbf{vladimir.shelomovskii@gmail.com, vvsss}

 





\textbf{Solution 6}


                    \begin{figure}[H]
                        \centering
                        \includegraphics[width=0.5\textwidth]{Solution/AIME_I_2014_15_Solution_Figure_3.png}
                    \end{figure}
                    The main element of the solution is the proof that $G$ is midpoint of $AC.$

 As in Solution 5 we get $\angle GED = \angle DBG =\gamma \implies$

 $\triangle BCG$ is isosceles triangle with $BG=CG.$

 Similarly $BG = AG \implies AG = CG = BG = \frac {AC}{2} =\frac {5}{2}.$

 
 \[\overset{\Large\frown} {FG} = 90^\circ – \overset{\Large\frown} {GD} =  90^\circ – 2\gamma \implies\]
 \[\overset{\Large\frown} {BFG} = 4\gamma + 90^\circ – 2\gamma = 90^\circ + 2\gamma \implies\]
 \[\angle BOG = 90^\circ + 2\gamma \implies \angle BGO = \angle GBO = 45^\circ - \gamma.\] Let $\hspace{10mm} BO = EO = DO = r \implies$
 \[BG = 2 r \cos(45^\circ – \gamma) = 2 r (\sin \gamma + \cos \gamma)\frac {\sqrt {2}}{2} =\]
 \[r \biggl(\frac {3}{5} + \frac {4}{5}\biggr) \sqrt {2} = r \frac {7 \sqrt{2}}{5} = \frac {5}{2}\implies\]
 \[r = \frac {25 \cdot \sqrt{2}}{28}\implies ED = \frac {25 \cdot \sqrt{2}}{14}\implies  \boxed{\textbf{041}}.\]\textbf{vladimir.shelomovskii@gmail.com, vvsss}

 





\textbf{Solution 7}

Let $(BEFGD) = \omega$.By Incenter-Excenter(Fact $5$), $F$ is the angle bisector of $\angle B$.Then by Ratio Lemma we have
 \[\frac{AG}{CG} = \frac{\sin(ABG)}{\sin(CBG)} \cdot \frac{AB}{BC} = \frac{\sin(GDE)}{\sin(DEG)} \cdot \frac{3}{4} = 1\]Thus, $G$ is the midpoint of $AC$.

 We can calculate $AF$ and $CF$ to be $\frac{15}{7}$ and $\frac{20}{7}$ respectively.And then by Power of a Point, we have$\newline$
 \[\operatorname{Pow}_{\omega}(A) = AE \cdot AB = AF \cdot AG \implies AE = \frac{25}{14}\]And then similarly, we have $CD = AE = \frac{25}{14}$.$\newline$

 Then $EB = \frac{17}{14}$ and $DB = \frac{31}{14}$ and by Pythagorean we have $DE = \frac{25\sqrt{2}}{14}$, so our answer is $\boxed{\textbf{041}}.$

 





\textbf{Solution 8 (funny angle chase & trig)}

Since $\angle EBD$ is right, $DE$ is clearly the diameter. Let $EF=DF=x, ED=x\sqrt{2}, DG=\frac{3x\sqrt{2}}{5}, EG=\frac{4x\sqrt{2}}{5}$. Then, let $\angle DEG=\alpha$. Therefore, $\angle EFG=\angle EFD+\angle DFG=90^{\circ}+\alpha,$ and $\angle EFA=90^{\circ}-\alpha$. However, $\triangle DGE \sim \triangle ABC$ so $\angle AFE$ also equals $90^{\circ}-\alpha$. Thus, $\triangle AEF$ is isosceles with $\angle FAE \cong \angle EFA \implies AE=EF=x.$

 Furthermore, $\angle FEG = \angle FDG = \angle EDG-\angle EDF=90^{\circ}-\alpha-45^{\circ}=45^{\circ}-\alpha$. Also, $\angle AEF=2\alpha$ in triangle $\triangle AEF$, thus $\angle BED=135^{\circ}-2\alpha$ since $\angle AEB=180^{\circ}$. Using $\cos \alpha=\frac{4}{5}$, it's relatively easy to derive that $\cos (135^{\circ}-2\alpha)=\frac{17\sqrt{2}}{50}$. Since $\cos(135^{\circ}-2\alpha)=\frac{BE}{DE}$, we get that $BE=\frac{17x}{25}$. Finally, since $AE+BE=x+\frac{17x}{25}=3$, we solve for 
 \[x=\frac{25}{14} \implies DE=x\sqrt{2}=\frac{25\sqrt{2}}{14},\] so our desired answer is $\boxed{041}.$

 





\textbf{Solution 9 (Vectors)}


                    \begin{figure}[H]
                        \centering
                        \includegraphics[width=0.5\textwidth]{Solution/AIME_I_2014_15_Solution_Figure_4.png}
                    \end{figure}
                    $\angle{B} = 90^\circ ==>$ $ED$ is the diameter of the circle $==> \angle{EGD} = 90^\circ$

 Because $\frac{GD}{GE} = \frac{AB}{BC} = \frac{3}{4}, \Delta DGE \sim \Delta ABC$

 $\angle{DFC} = \angle{DEG} = \angle{C} ==> CD = DF = EF$

 We flip the triangle for easier calculation and let $C$ be the origin.

 
                    \begin{figure}[H]
                        \centering
                        \includegraphics[width=0.5\textwidth]{Solution/AIME_I_2014_15_Solution_Figure_5.png}
                    \end{figure}
                    Let $x = CD$

 $DF = x\cdot e^{i2\alpha}, EF = x\cdot e^{-(90^\circ-2\alpha)}$

 $E$ is on $AB$, which means that $Re(E)$ = 4.

 
 \[Re(x+x\cdot e^{i2\alpha}+x\cdot e^{-(90^\circ-2\alpha)}) = 4\]

 
 \[Re(x+x(\cos{2\alpha}+i\sin{2\alpha})+x(\cos{(2\alpha-90^\circ)}+i\sin{(2\alpha-90^\circ)})) = 4\]

 
 \[x(1+\cos{2\alpha}+\sin{2\alpha})=4\]

 $\cos{2\alpha} = \cos^2{\alpha} - \sin^2{\alpha} = \frac{7}{25}$

 $\sin{2\alpha} = 2\sin{\alpha}\cos{\alpha} = \frac{24}{25}$

 
 \[x\cdot \frac{56}{25} = 4\]

 
 \[x = \frac{25}{14}\]

 Therefore $DE=\sqrt{2}\cdot \frac{25}{14} = \frac{25\sqrt{2}}{14}$ 

 $25+2+14=\boxed{041}$

 





\textbf{Solution 10 (Very similar to Solution 1 but with an added LoS)}

As noted in Solution 1, $DE$ is in fact the diameter of the circle. The proof is trivial. Note that $\triangle ABC$ is a right triangle which means $\angle ABC = \angle EBD = 90^{\circ}$. Therefore $\triangle EBD$ is an inscribed right triangle of the circle which means its arc length is just double the inscribed angle(following the Inscribed Angle Theorem) meaning the circle is cut in half by $DE \implies DE$ is indeed the diameter of the circle.

 Similarly, as $DE$ is the diameter, inscribed angles $\angle EFD = \angle EGD = 90^{\circ}$. As $EF = DF$, it follows $\angle FED = \angle FDE = 45^{\circ} \implies \triangle EGD$ is isosceles. 

 Let $EF = DF = y$:

 $ED = y\sqrt{2}$

 Now as $\frac{DG}{EG} = \frac{3}{4} \implies DG = 3x, EG = 4x$

 Hence $\triangle EGD$ is a right triangle with legs of length $3x, 4x$ and hypotenuse of length $y\sqrt{2}$. Obviously, by $SAS$ similarity, $\triangle EGD \sim \triangle ABC \implies \frac{DG}{ED} = \frac{AB}{AC} \implies \frac{3x}{y\sqrt{2}} = \frac{3}{5} \implies y\sqrt{2} = 5x \implies y = \frac{5x}{\sqrt{2}}$.

 Now note that quadrilateral $EFGD$ is cyclic which yields opposite angles are supplementary. Hence

 $\angle FED + \angle FGD = 180^{\circ}$ and $\angle FGD + \angle DGC = 180^{\circ}$ following Angle Addition Postulate. Hence by substitution, $\angle FED = \angle DGC = 45^{\circ}$. 

 Let $\angle ACB = \theta$. Hence $\sin(\theta) = \frac{3}{5}$ following the $3-4-5$ right triangle $ABC$. Now $\angle ACB = \angle GCD = \theta$ so we can apply Law of Sines on $\triangle DGC$. 

 We already know $\angle DGC = 45^{\circ}$. Also $\angle GCD = \theta$ and $GD = 3x$. Let $DC = a$. Now we apply Law of Sines:

 $\frac{DC}{\sin(\angle DGC)} = \frac{GD}{\sin(\angle GCD)} \implies \frac{a}{\sin(45^{\circ})} = \frac{3x}{\sin(\theta)} \implies \frac{a}{\frac{\sqrt{2}}{2}} = \frac{3x}{\frac{3}{5}} \implies a = \frac{5x}{\sqrt{2}}$. This means $DC = \frac{5x}{\sqrt{2}}$. 

 Now based on how we found $\angle FED = \angle DGC = 45^{\circ}$ using a combination of Supplementary Angle Theorem for cyclic quadrilaterals and Angle Addition Postulate, we can conclude based on a similar logic that $\angle EDG = \angle EFA$. The proof is below:

 $\angle EDG + \angle EFG = 180^{\circ}$ and $\angle EFG + \angle EFA = 180^{\circ}$. By substitution, $\angle EDG = \angle EFA$. 

 Now let $\angle EDG = \alpha$. From the right triangle $\triangle EDG$, $\sin(\angle EDG) = \frac{EG}{ED} = \frac{4x}{5x} = \frac{4}{5}$. So $\sin(\alpha) = \frac{4}{5}$. 

 Now let $\angle BAC = \angle EAF = \beta$. From the right triangle $\triangle ABC$, $\sin(\angle BAC) = \frac{4}{5} \implies \sin(\angle EAF) = \frac{4}{5} \implies \sin(\beta) = \frac{4}{5}$. 

 Now we apply Law of Sines on $\triangle EAF$:

 $\frac{EF}{\sin(\angle EAF)} = \frac{AE}{\sin(\angle EFA)} \implies \frac{\frac{5x}{\sqrt{2}}}{\sin(\angle \beta)} = \frac{AE}{\sin(\alpha)} \implies \frac{\frac{5x}{\sqrt{2}}}{\frac{4}{5}} = \frac{AE}{\frac{4}{5}} \implies AE = \frac{5x}{\sqrt{2}}$.

 Hence $EB = AB - AE = 3 - \frac{5x}{\sqrt{2}}$ and $BD = CB - CD = 4 - \frac{5x}{\sqrt{2}}$. 

 Now since $\triangle EBD$ is a right triangle with legs of $EB, BD$ and hypotenuse $ED = 5x$ we solve

 $(3 - \frac{5x}{\sqrt{2}})^{2} + (4 - \frac{5x}{\sqrt{2}})^{2} = (5x)^{2}$

 $9 - 6 \cdot \frac{5x}{\sqrt{2}} + \frac{25x^{2}}{2} + 16 - 8 \cdot \frac{5x}{\sqrt{2}} + \frac{25x^{2}}{2} = 25x^{2}$

 $25 - \frac{30x}{\sqrt{2}} - \frac{40x}{\sqrt{2}} = 0$

 $25 = \frac{70x}{\sqrt{2}} = 35x\sqrt{2}$. 

 Now recall we're looking for $ED = 5x$ so we divide both sides by $7\sqrt{2}$

 $\frac{25}{7\sqrt{2}} = 5x \implies 5x = \frac{25\sqrt{2}}{14}$.

 Thus the answer is $25 + 2 + 14 = \boxed{41}$. 

 





\textbf{Solution 11 (Coordinate Geometry)}

We are given that points $E,B,D,G,F$ are cyclic and $DF=EF$. Hence, we have $\angle EBF=\angle DBF$. We establish a coordinate system with $B$ being the origin. Here is our setup:\begin{align*}A&=(0,3)\\B&=(0,0)\\C&=(4,0)\\D&=(x_1,0)\\E&=(0,y_1)\\F&=(x_2,y_2)\\G&=(x_3,y_3)\\AC&:y=-\frac{3}{4}x+3\\\omega&: \left(x-\frac{x_{1}}{2}\right)^2+\left(y-\frac{y_{1}}{2}\right)^2=r^2\end{align*}
Since $\angle EBF=\angle DBF$, $F$ lies on the line $y=x$. We find $F=\left(\frac{12}{7},\frac{12}{7}\right)$. Also, the radius $r=\sqrt{\left(\frac{x_1}{2}\right)^2+\left(\frac{y_1}{2}\right)^2}$. Thus, the equation of $\omega$ simplifies to
 \[x^2-x_1x+y^2-y_1y=0\]$F$ lies on the circle $\omega$, we have\begin{gather*}\frac{144}{49}-\frac{12}{7}x_{1}+\frac{144}{49}-\frac{12}{7}y_{1}=0\\x_{1}+y_{1}=\frac{24}{7}\tag{1}\end{gather*}Solving the equation of $AC$ and $\omega$, we find\begin{align*}x^2-x_1x+(3-\frac{3}{4}x)^2-y_1(3-\frac{3}{4}x)&=0\\\frac{25}{16}x^2-(x_1+\frac{9}{2}-\frac{3}{4}y_{1})x+9-3y_{1}&=0\end{align*}
By Vieta's formula,
 \[\frac{12}{7}x_{3}=\frac{48(3-y_1)}{25}\]Solving yields $x_{3}=\frac{28(3-y_1)}{25}$, and $G \text{ lies on }AC\implies y_{3}=3-\frac{3}{4}x_3$. Therefore,
 \[G=\left(\frac{28(3-y_1)}{25}, \frac{12+21y_{1}}{25}\right)\]We proceed by solving $\frac{DG^2}{EG^2}=\frac{9}{16}$. By $(1)$, we substitute $x_{1}=\frac{24}{7}-y_1$. By the distance formula,\begin{align*}DG^2&=\left(\frac{28(3-y_1)}{25}-(\frac{24}{7}-y_1)\right)^2+\left(\frac{12+21y_{1}}{25}\right)^2\\DG^2&=\left(-\frac{3(4+7y_{1})}{175}\right)^2+\left(\frac{3(4+7y_{1})}{25}\right)^2\\DG^2&=\frac{18}{1225}(4+7y_{1})^2\end{align*}
Computing $EG^2$:\begin{align*}EG^2&=\left(\frac{28(3-y_1)}{25}\right)^2+\left(\frac{12+21y_{1}}{25}-y_1\right)^2\\EG^2&=\left(\frac{28(3-y_1)}{25}\right)^2+\left(\frac{4(3-y_{1})}{25}\right)^2\\EG^2&=\frac{32}{25}(3-y_{1})^2\end{align*}
Therefore,
 \[\frac{DG^2}{EG^2}=\frac{\frac{18}{1225}(4+7y_{1})^2}{\frac{32}{25}(3-y_{1})^2}=\frac{9}{16\cdot 49}\frac{(4+7y_{1})^2}{(3-y_{1})^2}=\frac{9}{16}\]Taking the positive solution from $\left\lvert\frac{4+7y_{1}}{3-y_{1}}\right\rvert=7$, we have $y_{1}=\frac{17}{14}$ and $x_{1}=\frac{31}{14}$. Therefore,
 \[DE=\sqrt{x_1^2+y_1^2}=\sqrt{\left(\frac{31}{14}\right)^2+\left(\frac{17}{14}\right)^2}=\frac{25\sqrt{2}}{14}\]The requested sum is $25+2+14=\boxed{041}$.

 \textbf{Fast Video Solution} 




\end{enumerate}